\documentclass[journal=jcisd8,manuscript=suppinfo]{achemso}

\usepackage{amsmath,amssymb}
\usepackage{graphicx}
\usepackage{booktabs}
\usepackage{multirow}
\usepackage{siunitx}
\usepackage[dvipsnames]{xcolor}
\usepackage{enumitem}

% Margin relief
\usepackage{microtype}
\setlength{\emergencystretch}{3em}
\setlength{\tabcolsep}{4pt}
\sloppy

% Mirror the vector/matrix shorthand from the main paper (needed for the
% moved equation blocks in the new Detailed Model Architectures section).
\newcommand{\vect}[1]{\mathbf{#1}}
\newcommand{\matr}[1]{\mathbf{#1}}

\title{Supporting Information: When Do Simple Models Win? Machine Learning Architectures for UV Absorption Prediction}

\author{Umesh Arampath}
\email{uarampath@mwbioprocessing.com}
\affiliation{Department of Mathematics, University of Iowa, Iowa City, USA}
\alsoaffiliation{Midwest Bioprocessing Center, Illinois, USA}

\author{Bryant Pero}
\affiliation{Midwest Bioprocessing Center, Illinois, USA}

\author{David Stewart}
\affiliation{Department of Mathematics, University of Iowa, Iowa City, USA}

\author{David Demirjian}
\affiliation{Midwest Bioprocessing Center, Illinois, USA}

\begin{document}

\section{Dataset Overview}
\label{si:dataset_overview}

Figure~\ref{fig:si_dataset_overview} provides a three-panel overview: (a)~the $\lambda_{\max}$ distribution across the $18{,}415$ records (median $= 396$~nm; range $[162, 1100]$~nm, with the bulk in the UV/early-visible region 250--500~nm and a long tail of deep-coloured chromophores extending to the near-IR); (b)~the top-20 solvents by record count, on a log scale to accommodate the $\sim$30$\times$ dynamic range between the most common (DCM, MeCN, CHCl$_3$) and least common solvents in the top-20 set (note the $726$ unique canonical solvent SMILES in the full dataset; the long tail of $\sim$700 rarely-used solvents is not shown); and (c)~eight representative chromophore structures spanning the major class diversity in the dataset (hydroxycinnamic acids, linear acenes, push--pull aromatics, azo dyes, BF$_2$-dipyrromethene BODIPYs, push--pull coumarins, extended-quinoid indigo, and porphyrin macrocycles).

\begin{figure}[htbp]
\centering
\includegraphics[width=\textwidth]{results/figure_dataset_overview.png}
\caption{Overview of the v3 cleaned Joung+Beard primary dataset ($18{,}415$ records, $7{,}846$ unique chromophores, $726$ unique canonical solvent SMILES). (a) $\lambda_{\max}$ histogram (bin width $\approx 10$~nm; red dashed line marks the median $396$~nm). (b) Top-20 solvents ranked by record count (log~$x$-axis; raw counts annotated). (c) Eight representative chromophore structures spanning class diversity. All numbers reproducible from \texttt{make\_dataset\_overview.py} in the public repository.}
\label{fig:si_dataset_overview}
\end{figure}

This panel-(a)/panel-(c) view also clarifies the structural diversity that motivates the locality-hypothesis discussion in the main text: the dataset includes both molecules dominated by local conjugation (e.g., hydroxycinnamic acids, BODIPYs) and molecules where extended conjugation and through-bond donor--acceptor charge transfer span the entire scaffold (e.g., porphyrins, indigo, push--pull cyanines). The non-local subset analysis (Section~\ref{si:nonlocal_smarts}, summarised in the Model Comparison section of the main text) partitions the dataset along this axis.


\section{Per-Fold Cross-Validation Results}
\label{sec:perfold}

Tables~\ref{tab:perfold_rf}--\ref{tab:perfold_chemberta} present per-fold results for all five model families on the primary Joung+Beard dataset (stratified 5-fold CV, 64/16/20 split).

\begin{table}[htbp]
\centering
\caption{Per-fold RF + Morgan FP results (tuned, with solvent) on the Joung+Beard v3 cleaned dataset (18{,}415 records after Greenman--Song duplicate handling).}
\label{tab:perfold_rf}
\begin{tabular}{ccccc}
\toprule
\textbf{Fold} & \textbf{RMSE (nm)} & \textbf{MAE (nm)} & $\boldsymbol{R^2}$ & \textbf{Pearson $r$} \\
\midrule
0 & 30.33 & 15.26 & 0.920 & 0.961 \\
1 & 30.46 & 15.25 & 0.917 & 0.958 \\
2 & 30.24 & 15.15 & 0.919 & 0.960 \\
3 & 32.67 & 15.41 & 0.908 & 0.953 \\
4 & 33.81 & 15.79 & 0.903 & 0.951 \\
\midrule
\textbf{Mean $\pm$ Std} & \textbf{31.50 $\pm$ 1.47} & \textbf{15.37 $\pm$ 0.22} & \textbf{0.914 $\pm$ 0.007} & \textbf{0.957 $\pm$ 0.004} \\
\bottomrule
\end{tabular}
\end{table}

\begin{table}[htbp]
\centering
\caption{Per-fold BiGRU + Solvent results (v2).}
\label{tab:perfold_bigru}
\begin{tabular}{ccccc}
\toprule
\textbf{Fold} & \textbf{RMSE (nm)} & \textbf{MAE (nm)} & $\boldsymbol{R^2}$ & \textbf{Pearson $r$} \\
\midrule
0 & 35.23 & 20.02 & 0.888 & 0.942 \\
1 & 37.33 & 20.86 & 0.879 & 0.937 \\
2 & 34.96 & 20.77 & 0.893 & 0.945 \\
3 & 37.62 & 21.51 & 0.877 & 0.937 \\
4 & 37.12 & 20.32 & 0.881 & 0.939 \\
\midrule
\textbf{Mean $\pm$ Std} & \textbf{36.45 $\pm$ 1.12} & \textbf{20.70 $\pm$ 0.51} & \textbf{0.884 $\pm$ 0.006} & \textbf{0.940 $\pm$ 0.003} \\
\bottomrule
\end{tabular}
\end{table}

\begin{table}[htbp]
\centering
\caption{Per-fold BiGRU (tuned, 3L/256u) results on the Joung+Beard v3 cleaned dataset.}
\label{tab:perfold_bigru_tuned}
\begin{tabular}{ccccc}
\toprule
\textbf{Fold} & \textbf{RMSE (nm)} & \textbf{MAE (nm)} & $\boldsymbol{R^2}$ & \textbf{Pearson $r$} \\
\midrule
0 & 35.42 & 18.14 & 0.891 & 0.944 \\
1 & 35.74 & 18.95 & 0.886 & 0.942 \\
2 & 34.77 & 18.54 & 0.893 & 0.945 \\
3 & 36.90 & 18.56 & 0.883 & 0.940 \\
4 & 38.15 & 19.84 & 0.877 & 0.937 \\
\midrule
\textbf{Mean $\pm$ Std} & \textbf{36.20 $\pm$ 1.20} & \textbf{18.81 $\pm$ 0.58} & \textbf{0.886 $\pm$ 0.006} & \textbf{0.941 $\pm$ 0.003} \\
\bottomrule
\end{tabular}
\end{table}

\begin{table}[htbp]
\centering
\caption{Per-fold XGBoost + Morgan FP results (MSE loss) on the Joung+Beard v3 cleaned dataset.}
\label{tab:perfold_xgboost}
\begin{tabular}{ccccc}
\toprule
\textbf{Fold} & \textbf{RMSE (nm)} & \textbf{MAE (nm)} & $\boldsymbol{R^2}$ & \textbf{Pearson $r$} \\
\midrule
0 & 33.41 & 20.27 & 0.903 & 0.952 \\
1 & 32.07 & 19.55 & 0.908 & 0.953 \\
2 & 33.59 & 20.64 & 0.900 & 0.950 \\
3 & 34.76 & 20.15 & 0.896 & 0.947 \\
4 & 35.77 & 20.66 & 0.892 & 0.945 \\
\midrule
\textbf{Mean $\pm$ Std} & \textbf{33.92 $\pm$ 1.26} & \textbf{20.25 $\pm$ 0.41} & \textbf{0.900 $\pm$ 0.006} & \textbf{0.949 $\pm$ 0.003} \\
\bottomrule
\end{tabular}
\end{table}

\begin{table}[htbp]
\centering
\caption{Per-fold Chemprop D-MPNN results on the Joung+Beard v3 cleaned dataset (multicomponent: solute + solvent).}
\label{tab:perfold_chemprop}
\begin{tabular}{ccccc}
\toprule
\textbf{Fold} & \textbf{RMSE (nm)} & \textbf{MAE (nm)} & $\boldsymbol{R^2}$ & \textbf{Pearson $r$} \\
\midrule
0 & 30.62 & 16.22 & 0.919 & 0.960 \\
1 & 30.53 & 16.72 & 0.917 & 0.958 \\
2 & 31.49 & 15.75 & 0.912 & 0.956 \\
3 & 33.90 & 17.87 & 0.901 & 0.950 \\
4 & 39.21 & 23.98 & 0.870 & 0.934 \\
\midrule
\textbf{Mean $\pm$ Std} & \textbf{33.15 $\pm$ 3.27} & \textbf{18.11 $\pm$ 3.02} & \textbf{0.904 $\pm$ 0.018} & \textbf{0.951 $\pm$ 0.009} \\
\bottomrule
\end{tabular}
\end{table}

\begin{table}[htbp]
\centering
\caption{Per-fold Chemprop D-MPNN results (v2, solute only, no solvent).}
\label{tab:perfold_chemprop_nosolvent}
\begin{tabular}{ccccc}
\toprule
\textbf{Fold} & \textbf{RMSE (nm)} & \textbf{MAE (nm)} & $\boldsymbol{R^2}$ & \textbf{Pearson $r$} \\
\midrule
0 & 33.37 & 15.95 & 0.899 & 0.949 \\
1 & 37.28 & 16.29 & 0.879 & 0.938 \\
2 & 29.31 & 15.88 & 0.925 & 0.962 \\
3 & 33.07 & 16.64 & 0.905 & 0.952 \\
4 & 35.88 & 17.64 & 0.889 & 0.943 \\
\midrule
\textbf{Mean $\pm$ Std} & \textbf{33.78 $\pm$ 2.73} & \textbf{16.48 $\pm$ 0.64} & \textbf{0.900 $\pm$ 0.015} & \textbf{0.949 $\pm$ 0.008} \\
\bottomrule
\end{tabular}
\end{table}

\begin{table}[htbp]
\centering
\caption{Per-fold ChemBERTa results on the Joung+Beard v3 cleaned dataset.}
\label{tab:perfold_chemberta}
\begin{tabular}{ccccc}
\toprule
\textbf{Fold} & \textbf{RMSE (nm)} & \textbf{MAE (nm)} & $\boldsymbol{R^2}$ & \textbf{Pearson $r$} \\
\midrule
0 & 50.03 & 23.98 & 0.783 & 0.905 \\
1 & 51.69 & 26.05 & 0.760 & 0.888 \\
2 & 56.09 & 28.29 & 0.722 & 0.875 \\
3 & 52.61 & 25.06 & 0.761 & 0.890 \\
4 & 60.02 & 30.52 & 0.696 & 0.868 \\
\midrule
\textbf{Mean $\pm$ Std} & \textbf{54.09 $\pm$ 3.57} & \textbf{26.78 $\pm$ 2.35} & \textbf{0.744 $\pm$ 0.031} & \textbf{0.885 $\pm$ 0.013} \\
\bottomrule
\end{tabular}
\end{table}


\section{Statistical Significance Tests}
\label{sec:significance}

Table~\ref{tab:significance} presents paired $t$-test results across the five cross-validation folds, computed on the v3 cleaned Joung+Beard dataset (18{,}415 records after Greenman--Song duplicate handling).
Most model comparisons are statistically significant at $\alpha = 0.05$ using the parametric paired $t$-test, with the notable exception of RF vs.\ Chemprop ($p = 0.16$ on v3, vs.\ $p = 0.77$ on the earlier v2 first-occurrence-dedup data: in both cases the null hypothesis of equal mean RMSE cannot be rejected, but the v3 $p$-value is closer to threshold because Chemprop fold~4 happens to be a high-variance early-stopping outlier on v3 at RMSE~$=$~39.21~nm).
With only $n = 5$ folds, the Wilcoxon signed-rank test (a non-parametric alternative) yields $p = 0.0625$ for all comparisons where the 5 paired differences are concordant (the minimum achievable $p$-value), providing strong directional evidence despite not reaching conventional significance.

\begin{table}[htbp]
\centering
\caption{Paired $t$-test results for RMSE comparisons across 5 v3 folds. Solvent vs.\ no-solvent rows are reported on the v2 dataset.}
\label{tab:significance}
\begin{tabular}{lccc}
\toprule
\textbf{Comparison} & $\boldsymbol{t}$ & $\boldsymbol{p}$ & \textbf{Significant?} \\
\midrule
RF vs.\ Chemprop & $-1.70$ & 0.164 & No \\
RF vs.\ BiGRU & $-22.54$ & $<$0.001 & Yes \\
RF vs.\ XGBoost & $-7.17$ & 0.002 & Yes \\
RF vs.\ ChemBERTa & $-15.78$ & $<$0.001 & Yes \\
Chemprop vs.\ XGBoost & $-0.70$ & 0.522 & No \\
Chemprop vs.\ BiGRU & $-2.74$ & 0.052 & Borderline \\
Chemprop vs.\ ChemBERTa & $-20.54$ & $<$0.001 & Yes \\
XGBoost vs.\ BiGRU & $-5.65$ & 0.005 & Yes \\
XGBoost vs.\ ChemBERTa & $-14.20$ & $<$0.001 & Yes \\
BiGRU vs.\ ChemBERTa & $-11.69$ & $<$0.001 & Yes \\
\midrule
RF: solvent vs.\ no-solvent (v2) & $-11.46$ & $<$0.001 & Yes \\
Chemprop: solvent vs.\ no-solvent (v2) & $-4.33$ & 0.012 & Yes \\
BiGRU: solvent vs.\ no-solvent (v2) & $-3.90$ & 0.018 & Yes \\
\bottomrule
\end{tabular}
\end{table}


\section{RF Hyperparameter Sensitivity}
\label{sec:rf_tuning}

The RF hyperparameter grid search evaluated 432 configurations varying the number of trees ($B \in \{100, 500, 1000\}$), maximum features per split ($\sqrt{d}$, $0.1d$, $0.3d$), maximum tree depth (None, 30, 50), minimum samples per leaf (1, 2, 5), fingerprint radius (2, 3), and fingerprint length (1024, 2048).
The search was conducted on fold~0 using the validation set for model selection.

The most influential parameter was \texttt{max\_features}: all top-10 configurations used $0.3d$ (30\% of features per split), compared to the scikit-learn default of $\sqrt{d}$ ($\approx$1.6\% of 4096 features).
This finding suggests that Morgan fingerprint features are sufficiently decorrelated that broader feature sampling per split improves performance.
The number of trees ($B$) showed diminishing returns beyond 500, and maximum depth had negligible effect (unlimited depth was optimal), indicating that the dataset does not suffer from overfitting with deep trees.


\section{BiGRU Training Curves}
\label{sec:learning_curves}

Figure~\ref{fig:learning_curves} shows the training and validation MAE curves for the BiGRU models, averaged across folds with $\pm 1$ standard deviation bands.
Both models converge within 75--125 epochs, with early stopping preventing overfitting.
The gap between training and validation curves indicates mild overfitting, consistent with the relatively small dataset size (12K training samples per fold).

\begin{figure}[htbp]
 \centering
 \includegraphics[width=0.9\textwidth]{results/learning_curves.png}
 \caption{BiGRU training curves showing training and validation MAE vs.\ epoch, averaged across folds. Shaded regions denote $\pm 1$ standard deviation. Left: BiGRU with solvent encoding. Right: BiGRU with solute only.}
 \label{fig:learning_curves}
\end{figure}


\section{ChemBERTa Learning Curves}
\label{sec:chemberta_curves}

Figure~\ref{fig:chemberta_curves} shows the ChemBERTa training and validation MAE curves across all five folds.
All folds converge within 40--60 epochs, with the final validation MAE plateauing around 24--26~nm.
The gap between training and validation MAE is larger than for BiGRU, consistent with ChemBERTa's 44.1M parameters being overparameterized for the $\sim$12K training samples available per fold.

\begin{figure}[htbp]
 \centering
 \includegraphics[width=\textwidth]{results/chemberta_learning_curves.png}
 \caption{ChemBERTa fine-tuning learning curves. Left: per-fold validation MAE (solid) and training MAE (dashed). Right: mean $\pm$ 1 standard deviation across folds, with final values annotated.}
 \label{fig:chemberta_curves}
\end{figure}


\section{Hyperparameter Tuning Asymmetry}
\label{sec:tuning_asymmetry}

Table~\ref{tab:tuning_asymmetry} compares the effect of hyperparameter tuning for both RF and BiGRU, demonstrating that the in-distribution performance gap between model families is not an artifact of asymmetric tuning effort.

\begin{table}[htbp]
\centering
\caption{Effect of hyperparameter tuning on cross-validated performance (stratified 5-fold CV).}
\label{tab:tuning_asymmetry}
\begin{tabular}{lccc}
\toprule
\textbf{Model} & \textbf{RMSE (nm)} & \textbf{MAE (nm)} & \textbf{Improvement} \\
\midrule
RF (default, $\sqrt{d}$) & 32.18 $\pm$ 1.74 & 15.42 $\pm$ 0.47 & --- \\
RF (tuned, $0.3d$) & 31.34 $\pm$ 1.82 & 15.16 $\pm$ 0.44 & $-$2.6\% RMSE \\
\midrule
BiGRU (default, 2L/128u) & 36.45 $\pm$ 1.12 & 20.70 $\pm$ 0.51 & --- \\
BiGRU (tuned, 3L/256u) & 34.71 $\pm$ 1.40 & 18.09 $\pm$ 0.59 & $-$4.8\% RMSE \\
\bottomrule
\end{tabular}
\end{table}

BiGRU HPO was conducted using Optuna TPE (25 trials on fold~0), with the best configuration extended to full 5-fold cross-validation, yielding a 4.8\% RMSE improvement over the default architecture.
Even the default RF (RMSE = 32.18~nm) outperforms the tuned BiGRU (RMSE = 34.71~nm), while RF tuning provides only an additional 2.6\% improvement.
BiGRU tuning narrows the gap from 14\% to 10\%, but the remaining difference suggests that the fingerprint representation has an inherent advantage for in-distribution UV absorption prediction.
ChemBERTa hyperparameters were not tuned due to the prohibitive training cost ($\sim$3 hours per fold for 44.1M parameters).


\section{Non-Local Subset Analysis: SMARTS Rules, Subset Definitions, and Per-Model Metrics}
\label{si:nonlocal_smarts}

This section reports the full non-local subset analysis that the Model Comparison section of the main text summarises in a two-sentence pointer. The primary benchmarks pool molecules whose absorption is plausibly governed by local substitution patterns with molecules where the chromophore spans the entire $\pi$-system, donor--acceptor charge-transfer dyes, low-gap polymethines, BODIPYs, and other classes where global electronic effects and longer conjugation paths are expected to dominate. To examine whether the locality argument continues to hold for this second class, we partitioned the Joung+Beard test folds into subsets enriched for non-local effects and recomputed per-model error metrics within each subset.

Subset definitions: \textit{long-wavelength} compounds ($\lambda_{\max} > 500$~nm or $> 600$~nm); \textit{donor--acceptor} (D-A) molecules containing at least one $\pi$-donor substituent (aryl-NH$_2$, aryl-NHR, aryl-NR$_2$, aryl-OH, aryl-OR, aryl-SH, aryl-SR) \emph{and} at least one $\pi$-acceptor (NO$_2$, CN, CHO, ketone, COOH/COOR/amide, SO$_2$R, CF$_3$, pyridinium, dicyanovinyl) attached to a conjugated $sp^2$ centre; \textit{D-A-D} molecules with $\geq 2$ donor matches and $\geq 1$ acceptor; \textit{azo} dyes (Ar--N$=$N--Ar); and \textit{BODIPY} cores (BF$_2$ bridging two ring nitrogens). The exact SMARTS patterns are listed in Tables~\ref{tab:si_donor_smarts}--\ref{tab:si_motif_smarts}, per-rule molecule counts in Table~\ref{tab:si_subset_counts}, and the full per-model error breakdown across all subsets in Table~\ref{tab:si_nonlocal_full}. The union of all non-local categories covers 49\% of the dataset (9{,}057 / 18{,}415 molecules), with the individual subset sizes ranging from 264 (azo) to 5{,}945 (D-A) molecules.

\textit{Donors} are lone-pair-bearing substituents on an aromatic carbon that push electron density into the $\pi$ system. The amide nitrogen is explicitly excluded (\texttt{!\$(N-C=O)}) because the acyl group converts the nitrogen into a weak acceptor rather than a donor.
\textit{Acceptors} are strongly electron-withdrawing substituents on a conjugated $sp^2$ centre---aromatic carbon (\texttt{[c]}) or an $sp^2$ alkene carbon (\texttt{[CX3](=[\#6,\#7])[!\#1]}). The alkene branch is included so that push--pull cinnamic acids such as ferulic acid (\texttt{COc1cc(/C=C/C(=O)O)ccc1O}) are correctly captured as donor--acceptor compounds.
The donor-count column shows the number of distinct atom-set matches per molecule (using RDKit's \texttt{GetSubstructMatches} with \texttt{uniquify=True}); a molecule with two methoxy groups on the same ring contributes a donor count of 2 to the D-A-D rule.

\begin{table}[htbp]
\centering
\caption{Donor SMARTS patterns used to identify D-A and D-A-D molecules.}
\label{tab:si_donor_smarts}
\scriptsize
\begin{tabular}{lll}
\toprule
\textbf{Rule name} & \textbf{SMARTS} & \textbf{Matches} \\
\midrule
aryl\_NH2 & \verb|[c]-[NX3;H2;!$(N(C=O));!$(N-N=O)]| & primary aryl amine \\
aryl\_NHR & \verb|[c]-[NX3;H1;!$(N-C=O);!$(N-N=O);!$(N-S(=O))]| & secondary aryl amine (non-amide) \\
aryl\_NR2 & \verb|[c]-[NX3;H0;!$(N-C=O);!$(N-N=O);!$(N-S(=O))]([#6])[#6]| & tertiary aryl amine \\
aryl\_OH & \verb|[c]-[OX2H1]| & aryl hydroxyl \\
aryl\_OR & \verb|[c]-[OX2;H0]-[CX4]| & aryl ether/methoxy \\
aryl\_SH & \verb|[c]-[SX2H1]| & aryl thiol \\
aryl\_SR & \verb|[c]-[SX2;H0]-[CX4]| & aryl sulfide \\
\bottomrule
\end{tabular}
\end{table}

\begin{table}[htbp]
\centering
\caption{Acceptor SMARTS patterns. The \emph{sp$^2$-anchor} \texttt{[\$([c]),\$([CX3](=[\#6,\#7])[!\#1])]} matches either an aromatic carbon or an alkene/imine sp$^2$ carbon so that push--pull cinnamic-acid-type acceptors are detected.}
\label{tab:si_acceptor_smarts}
\small
\begin{tabular}{lll}
\toprule
\textbf{Rule name} & \textbf{SMARTS} & \textbf{Matches} \\
\midrule
nitro & \verb|sp2-[$([NX3+](=O)[O-]),$([NX3](=O)=O)]| & nitro \\
CN & \verb|sp2-[CX2]#[NX1]| & cyano / nitrile \\
CHO & \verb|sp2-[CX3H1](=O)| & aldehyde \\
ketone & \verb|sp2-[CX3](=O)[#6]| & ketone \\
COOH & \verb|sp2-[CX3](=O)[OX2H1]| & carboxylic acid \\
COOR & \verb|sp2-[CX3](=O)[OX2H0][#6]| & ester \\
amide & \verb|sp2-[CX3](=O)[NX3]| & amide \\
SO2 & \verb|sp2-[SX4](=O)(=O)| & sulfonyl \\
CF3 & \verb|sp2-[CX4](F)(F)F| & trifluoromethyl \\
aromatic\_N+ & \verb|[c][n+,N+;X3]| & pyridinium / acridinium \\
dicyanovinyl & \verb|[#6]=C([CX2]#[NX1])[CX2]#[NX1]| & dicyanovinyl \\
malononitrile & \verb|[CX4H1]([CX2]#N)[CX2]#N| & malononitrile \\
\bottomrule
\end{tabular}
\end{table}

\begin{table}[htbp]
\centering
\caption{Specific dye motifs.}
\label{tab:si_motif_smarts}
\scriptsize
\begin{tabular}{lll}
\toprule
\textbf{Rule name} & \textbf{SMARTS} & \textbf{Matches} \\
\midrule
azo & \verb|[c]-[N;X2]=[N;X2]-[c]| & aryl azo bridge \\
bodipy & \verb|[B](F)(F)([#7;R])[#7;R]| & BF$_2$ bonded to two ring nitrogens \\
cyanine & \verb|[NX3,nX3]-,=[CX3]=,-[CX3]-,=[CX3]=,-[CX3]-,=[NX3+,nX3+]| & open-chain polymethine cyanine \\
squaraine & \verb|C1(=O)C(=*)C(=O)C1=*| & squaric-acid-derived push--pull \\
stilbene & \verb|[c]-[CH]=[CH]-[c]| & 1,2-diarylethene scaffold \\
\bottomrule
\end{tabular}
\end{table}

\begin{table}[htbp]
\centering
\caption{Non-local subset sizes in the full v3 cleaned Joung+Beard primary dataset ($n_{\text{total}} = 18{,}415$). Subset definitions are unioned with each test fold to obtain the per-fold $n$ values reported in Table~\ref{tab:si_nonlocal_full}.}
\label{tab:si_subset_counts}
\small
\begin{tabular}{lrr}
\toprule
\textbf{Subset} & $n$ & \textbf{\% of dataset} \\
\midrule
Full dataset & 18{,}415 & 100.00 \\
$\lambda_{\max} > 500$~nm & 3{,}832 & 20.81 \\
$\lambda_{\max} > 600$~nm & 1{,}555 & 8.44 \\
Azo (Ar--N$=$N--Ar) & 264 & 1.43 \\
D-A & 5{,}945 & 32.28 \\
D-A-D & 2{,}181 & 11.84 \\
BODIPY core & 1{,}714 & 9.31 \\
Cyanine & 59 & 0.32 \\
Stilbene & 1{,}512 & 8.21 \\
Squaraine & 0 & 0.00 \\
\midrule
Union (any non-local) & 9{,}057 & 49.18 \\
\bottomrule
\end{tabular}
\end{table}

\begin{table}[htbp]
\centering
\caption{Full per-model error breakdown on the non-local subsets of the v3 cleaned Joung+Beard dataset (mean $\pm$ std over the five stratified test folds). Best result per metric per subset in bold. The Squaraine subset is empty in the Joung+Beard dataset and is omitted. Reproduced via \texttt{python3 analyze\_nonlocal\_subsets.py --v3 --regen\_xgboost}.}
\label{tab:si_nonlocal_full}
\scriptsize
\setlength{\tabcolsep}{3pt}
\begin{tabular}{lcccccc}
\toprule
\textbf{Subset / metric} & $n$ & \textbf{RF} & \textbf{Chemprop} & \textbf{BiGRU} & \textbf{ChemBERTa} & \textbf{XGBoost} \\
\midrule
\multicolumn{7}{l}{\textbf{RMSE (nm)}} \\
Full dataset & 3683 & \textbf{31.50 $\pm$ 1.64} & 33.15 $\pm$ 3.65 & 36.20 $\pm$ 1.34 & 54.09 $\pm$ 3.99 & 33.92 $\pm$ 1.40 \\
$\lambda > 500$ & 766 & 51.81 $\pm$ 3.91 & \textbf{47.02 $\pm$ 7.78} & 57.45 $\pm$ 3.96 & 105.81 $\pm$ 8.46 & 51.20 $\pm$ 3.60 \\
$\lambda > 600$ & 311 & 69.54 $\pm$ 4.20 & \textbf{63.19 $\pm$ 8.67} & 77.14 $\pm$ 4.93 & 159.06 $\pm$ 12.48 & 67.63 $\pm$ 5.11 \\
Azo & 53 & \textbf{33.84 $\pm$ 7.40} & 37.36 $\pm$ 9.22 & 38.86 $\pm$ 8.26 & 45.62 $\pm$ 7.07 & 35.61 $\pm$ 3.64 \\
D-A & 1189 & 25.04 $\pm$ 0.41 & \textbf{24.54 $\pm$ 4.77} & 27.78 $\pm$ 1.31 & 47.26 $\pm$ 3.34 & 27.21 $\pm$ 1.11 \\
D-A-D & 436 & 26.86 $\pm$ 2.08 & \textbf{25.68 $\pm$ 5.83} & 29.25 $\pm$ 1.92 & 42.59 $\pm$ 4.29 & 28.50 $\pm$ 1.33 \\
BODIPY & 343 & \textbf{21.59 $\pm$ 2.14} & 22.83 $\pm$ 5.03 & 28.79 $\pm$ 0.64 & 78.02 $\pm$ 8.50 & 24.47 $\pm$ 2.69 \\
Cyanine & 12 & \textbf{17.39 $\pm$ 3.24} & 30.07 $\pm$ 4.12 & 22.15 $\pm$ 5.19 & 42.59 $\pm$ 10.75 & 25.55 $\pm$ 6.99 \\
Stilbene & 302 & \textbf{29.98 $\pm$ 1.54} & 33.36 $\pm$ 5.91 & 35.57 $\pm$ 5.28 & 50.40 $\pm$ 6.55 & 34.97 $\pm$ 3.32 \\
Non-local union & 1811 & 36.76 $\pm$ 2.38 & \textbf{35.03 $\pm$ 5.66} & 41.54 $\pm$ 2.41 & 71.11 $\pm$ 5.81 & 37.80 $\pm$ 2.13 \\
\midrule
\multicolumn{7}{l}{\textbf{MAE (nm)}} \\
Full dataset & 3683 & \textbf{15.37 $\pm$ 0.25} & 18.11 $\pm$ 3.38 & 18.81 $\pm$ 0.64 & 26.78 $\pm$ 2.63 & 20.25 $\pm$ 0.45 \\
$\lambda > 500$ & 766 & 25.38 $\pm$ 1.66 & \textbf{24.66 $\pm$ 6.28} & 29.20 $\pm$ 2.48 & 72.76 $\pm$ 8.82 & 28.42 $\pm$ 1.28 \\
$\lambda > 600$ & 311 & 35.61 $\pm$ 2.66 & \textbf{31.62 $\pm$ 6.64} & 40.25 $\pm$ 2.22 & 141.62 $\pm$ 13.94 & 37.02 $\pm$ 2.01 \\
Azo & 53 & \textbf{21.06 $\pm$ 3.70} & 25.39 $\pm$ 6.79 & 25.27 $\pm$ 4.38 & 26.66 $\pm$ 4.75 & 25.04 $\pm$ 2.00 \\
D-A & 1189 & \textbf{12.92 $\pm$ 0.25} & 15.73 $\pm$ 3.70 & 15.59 $\pm$ 0.76 & 23.19 $\pm$ 2.19 & 17.54 $\pm$ 0.66 \\
D-A-D & 436 & \textbf{14.38 $\pm$ 0.85} & 16.56 $\pm$ 3.84 & 17.11 $\pm$ 0.90 & 22.68 $\pm$ 2.25 & 18.58 $\pm$ 0.44 \\
BODIPY & 343 & \textbf{10.21 $\pm$ 0.75} & 15.32 $\pm$ 4.40 & 15.75 $\pm$ 1.41 & 50.70 $\pm$ 7.82 & 15.18 $\pm$ 0.99 \\
Cyanine & 12 & \textbf{14.43 $\pm$ 3.03} & 23.71 $\pm$ 3.34 & 18.18 $\pm$ 5.47 & 35.71 $\pm$ 10.55 & 22.25 $\pm$ 6.16 \\
Stilbene & 302 & \textbf{16.33 $\pm$ 1.34} & 20.79 $\pm$ 3.87 & 19.21 $\pm$ 1.92 & 27.91 $\pm$ 3.92 & 22.48 $\pm$ 1.35 \\
Non-local union & 1811 & \textbf{16.95 $\pm$ 0.75} & 18.89 $\pm$ 4.43 & 20.44 $\pm$ 1.34 & 38.22 $\pm$ 4.50 & 21.21 $\pm$ 0.58 \\
\midrule
\multicolumn{7}{l}{$\boldsymbol{R^2}$} \\
Full dataset & 3683 & \textbf{0.914 $\pm$ 0.007} & 0.904 $\pm$ 0.020 & 0.886 $\pm$ 0.007 & 0.744 $\pm$ 0.035 & 0.900 $\pm$ 0.006 \\
$\lambda > 500$ & 766 & 0.599 $\pm$ 0.062 & \textbf{0.665 $\pm$ 0.115} & 0.507 $\pm$ 0.073 & $-0.675 \pm 0.285$ & 0.610 $\pm$ 0.048 \\
$\lambda > 600$ & 311 & $-0.140 \pm 0.183$ & \textbf{0.049 $\pm$ 0.285} & $-0.394 \pm 0.142$ & $-4.967 \pm 1.071$ & $-0.072 \pm 0.140$ \\
Azo & 53 & \textbf{0.682 $\pm$ 0.106} & 0.617 $\pm$ 0.123 & 0.566 $\pm$ 0.194 & 0.424 $\pm$ 0.126 & 0.650 $\pm$ 0.042 \\
D-A & 1189 & 0.934 $\pm$ 0.004 & \textbf{0.935 $\pm$ 0.026} & 0.919 $\pm$ 0.005 & 0.764 $\pm$ 0.032 & 0.922 $\pm$ 0.007 \\
D-A-D & 436 & 0.916 $\pm$ 0.013 & \textbf{0.921 $\pm$ 0.038} & 0.901 $\pm$ 0.010 & 0.788 $\pm$ 0.044 & 0.906 $\pm$ 0.010 \\
BODIPY & 343 & \textbf{0.940 $\pm$ 0.009} & 0.930 $\pm$ 0.032 & 0.893 $\pm$ 0.006 & 0.208 $\pm$ 0.167 & 0.923 $\pm$ 0.013 \\
Cyanine & 12 & \textbf{0.856 $\pm$ 0.063} & 0.576 $\pm$ 0.135 & 0.768 $\pm$ 0.109 & 0.206 $\pm$ 0.205 & 0.679 $\pm$ 0.156 \\
Stilbene & 302 & \textbf{0.903 $\pm$ 0.012} & 0.877 $\pm$ 0.043 & 0.862 $\pm$ 0.040 & 0.723 $\pm$ 0.072 & 0.868 $\pm$ 0.021 \\
Non-local union & 1811 & 0.904 $\pm$ 0.012 & \textbf{0.912 $\pm$ 0.030} & 0.878 $\pm$ 0.014 & 0.641 $\pm$ 0.060 & 0.899 $\pm$ 0.010 \\
\bottomrule
\end{tabular}
\end{table}

\paragraph{Two trends across the subset table.}

\textit{First}, every model degrades sharply on long-wavelength compounds: RMSE roughly doubles when restricting to $\lambda_{\max} > 500$~nm and more than doubles for $\lambda_{\max} > 600$~nm. The effect is seen across all five architectures, but ChemBERTa exhibits an extreme version of it (RMSE = 106 and 159~nm respectively, with negative $R^2$ on both subsets): its predictions collapse to an essentially constant value around 600~nm, the prediction-ceiling artefact visible in the main-text parity plot and diagnosed in detail in the Error Analysis section of the main text. For the other four models the degradation is consistent with two compounding factors identified in prior work on these datasets~\cite{beard2019comparative,joung2021deep}: data scarcity in the deep-coloured region of $\lambda$-space, and the greater structural complexity of low-gap chromophores (extended polyenes, push--pull cyanines, BODIPYs, phthalocyanines), many of which carry charge-transfer character whose accurate description requires global rather than local features.

\textit{Second}, the rank between RF and Chemprop is \emph{not stable} across subsets, but the direction of the instability is informative. On the full dataset the two are statistically indistinguishable ($p = 0.16$; see the Model Comparison section of the main text), but on the long-wavelength subsets Chemprop pulls clearly ahead of RF: by 4.8~nm on $\lambda_{\max} > 500$~nm and 6.4~nm on $\lambda_{\max} > 600$~nm. On the broader D-A and D-A-D subsets Chemprop also retains a small advantage (0.5~nm on D-A, 1.2~nm on D-A-D), but those margins are smaller than they were on the v2 first-occurrence-dedup data; the qualitative direction (Chemprop $\geq$ RF on D-A-style compounds) is preserved but the v2-to-v3 protocol change has narrowed the gap, partly because Chemprop's fold-to-fold variance widened on v3 (std~$=$~4.77 on D-A, 5.83 on D-A-D). On two subsets the pattern flips: RF beats Chemprop on the small azo set (n~$=$~53; 33.8 vs.\ 37.4~nm, within one fold std of either model so not a robust reversal) and RF maintains an edge on BODIPYs (21.6 vs.\ 22.8~nm), a class whose rigid, conserved scaffold means that absorption is largely explained by short-range substitution patterns the fingerprint already captures. Taken together, the subset analysis is consistent with the locality argument (message-passing GNNs benefit relatively more when the property is governed by global rather than local features, with the largest effect seen on the deep-coloured tail) without overturning the headline observation that RF and Chemprop are within noise of each other on the dataset as a whole.

\paragraph{Additional observations.}
(i) The MAE column reinforces the RMSE story for the long-wavelength subsets ($\lambda > 500$ and $\lambda > 600$~nm: Chemprop best on both metrics) but flips on the local-context subsets (D-A, D-A-D, BODIPY, Stilbene, Cyanine, and now also Azo: RF lowest MAE despite being tied or behind on RMSE for some subsets). This is consistent with RF's residuals being concentrated near zero with a longer tail of larger errors, while Chemprop produces a more uniformly spread residual distribution---squared-error metrics weight the tail more heavily.
(ii) The cyanine subset contains only $\sim$12 test molecules per fold, so its metrics carry large fold-to-fold variance and should not be read as a reliable model ranking.
(iii) ChemBERTa's $R^2 = -4.97$ on the $\lambda > 600$~nm subset reflects the prediction-ceiling artefact discussed in the Error Analysis section of the main text---ChemBERTa is essentially predicting a constant in this region, so the variance-explained metric goes strongly negative.


\section{Overlap Between Primary and Deep4Chem Datasets}
\label{si:deep4chem_overlap}

We quantified the chemical-space overlap between our primary Joung+Beard dataset and the Deep4Chem external benchmark used in main-text the Cross-Dataset Benchmarking section. Joung et al.~\cite{joung2021deep} derived Deep4Chem from a similar literature curation effort to the Joung half of our primary collection, so a non-trivial overlap is expected; the open question is how large it actually is. All SMILES strings were canonicalised through RDKit~\cite{landrumRdkitRdkit2020_03_12020} before set intersection, separately for solutes and for solvents.

Table~\ref{tab:si_d4c_overlap} reports the overlap at three granularities. At the level of unique canonical solute SMILES, 77.4\% of Deep4Chem solutes (5{,}018 of 6{,}480) also appear in Joung+Beard; conversely, 63.8\% of Joung+Beard solutes (5{,}018 of 7{,}863) appear in Deep4Chem. The overlap is even more striking at the (solute, solvent) pair level---72.6\% of Deep4Chem records have an exact match in Joung+Beard---and is near-total when measured at the level of generic Bemis--Murcko scaffolds (98.8\%). The scaffold overlap shows that the two datasets explore essentially the same scaffold space, with Deep4Chem providing only modest scaffold-level novelty beyond Joung+Beard.

\begin{table}[htbp]
\centering
\caption{Overlap between the primary Joung+Beard dataset and the Deep4Chem external dataset, measured at three granularities after RDKit canonicalisation.}
\label{tab:si_d4c_overlap}
\small
\begin{tabular}{lrrrrr}
\toprule
\textbf{Granularity} & \textbf{J+B} & \textbf{D4C} & \textbf{Intersection} & \textbf{\% of D4C} & \textbf{\% of J+B} \\
\midrule
Unique canonical solute SMILES & 7{,}863 & 6{,}480 & 5{,}018 & 77.4 & 63.8 \\
(Solute, solvent) pairs & 18{,}457 & 17{,}427 & 12{,}654 & 72.6 & 68.6 \\
Generic Bemis--Murcko scaffolds & 3{,}008 & 2{,}498 & 2{,}467 & 98.8 & 82.0 \\
\bottomrule
\end{tabular}
\end{table}

\begin{figure}[htbp]
 \centering
 \includegraphics[width=0.95\textwidth]{results/r1_7_overlap_venn.png}
 \caption{Venn diagrams showing the chemical-space overlap between the Joung+Beard primary dataset and the Deep4Chem external benchmark. Left: unique canonical solute SMILES. Right: unique (solute, solvent) pairs after canonicalising both columns.}
 \label{fig:si_d4c_venn}
\end{figure}

The practical consequence for the main-text Cross-Dataset Benchmarking (the Cross-Dataset Benchmarking section) is that Deep4Chem results characterise \emph{in-distribution} interpolation rather than out-of-distribution extrapolation: roughly three-quarters of the Deep4Chem test set, under the published 80/10/10 random split, will contain chromophores the model has already learned from during training on Joung+Beard. We retain the Deep4Chem evaluation because it reproduces the published baselines of Joung et al.\ (GCNN, RMSE = 26.6~nm) on their own data, but we are explicit in the main text that the more rigorous external test is provided by Jung et al.~\cite{jung2024machine} (48.8\% unique-solute overlap, separate experimental provenance, distinct chromophore families). The scripts used to compute these overlaps are distributed with the code (\texttt{analyze\_deep4chem\_overlap.py}); the numbers above are reproducible from the source CSVs.


\section{v2-vs-v3 Cross-Protocol Comparison}
\label{si:v2_v3_comparison}

We adopted the Greenman--Song duplicate-handling protocol~\cite{greenman2022multifidelity,song2025fluortools} (drop within-key groups with $\max - \min > 5$~nm, average the rest, singletons unchanged) and trained all five model families on the resulting v3 cleaned dataset (18{,}415 records, vs.\ 18{,}755 under an alternative first-occurrence-dedup rule). Table~\ref{tab:si_v2v3} reports the side-by-side cross-validated metrics under both protocols.

The protocol change costs 1.81\% of records overall and re-distributes them across the deep-coloured and well-populated wavelength bins (Section~\ref{si:nonlocal_smarts}). For the four local-feature models the cross-validated RMSE shifts by $\leq 1.5$~nm; for ChemBERTa it shifts by $+3.7$~nm, with fold-to-fold standard deviation also widening (2.30 to 3.57). The headline RF~$\approx$~Chemprop tie is preserved: paired $t$-test $p = 0.16$ on v3 (vs.\ $p = 0.77$ on v2), so the null hypothesis of equal mean RMSE cannot be rejected at $\alpha = 0.05$ on either protocol. The qualitative model ranking---RF $\leq$ Chemprop $<$ XGBoost $<$ BiGRU $\ll$ ChemBERTa---is unchanged.

\begin{table}[htbp]
\centering
\caption{Cross-validated 5-fold metrics on the Joung+Beard primary dataset under the original first-occurrence-dedup protocol (v2, 18{,}755 records) and the Greenman--Song-style cleaned protocol adopted for the revised manuscript (v3, 18{,}415 records). $\Delta$ is the change in mean. Negative $\Delta R^2$ indicates a slight degradation under the cleaner protocol, which we attribute mainly to the test-fold composition shift rather than to the protocol per se. The BiGRU + Solvent (default) and XGBoost MAE rows are reported on v2 only.}
\label{tab:si_v2v3}
\small
\resizebox{\textwidth}{!}{%
\begin{tabular}{lcccccc}
\toprule
\textbf{Model} & \multicolumn{2}{c}{\textbf{RMSE (nm)}} & \multicolumn{2}{c}{\textbf{MAE (nm)}} & \multicolumn{2}{c}{$\boldsymbol{R^2}$} \\
\cmidrule(lr){2-3} \cmidrule(lr){4-5} \cmidrule(lr){6-7}
 & v2 & v3 & v2 & v3 & v2 & v3 \\
\midrule
RF + Morgan FP (tuned) & 31.34 $\pm$ 1.82 & \textbf{31.50 $\pm$ 1.47} & 15.16 $\pm$ 0.44 & 15.37 $\pm$ 0.22 & 0.914 & 0.914 \\
XGBoost + Morgan FP (MSE) & 33.70 $\pm$ 1.61 & 33.92 $\pm$ 1.26 & 20.05 $\pm$ 0.48 & 20.25 $\pm$ 0.41 & 0.900 & 0.900 \\
Chemprop (D-MPNN) & 31.69 $\pm$ 3.18 & 33.15 $\pm$ 3.27 & 16.84 $\pm$ 1.74 & 18.11 $\pm$ 3.02 & 0.911 & 0.904 \\
BiGRU (tuned, 3L/256u) & 34.71 $\pm$ 1.40 & 36.20 $\pm$ 1.20 & 18.09 $\pm$ 0.59 & 18.81 $\pm$ 0.58 & 0.894 & 0.886 \\
ChemBERTa (pretrained) & 50.39 $\pm$ 2.30 & 54.09 $\pm$ 3.57 & 24.31 $\pm$ 1.04 & 26.78 $\pm$ 2.35 & 0.777 & 0.744 \\
\bottomrule
\end{tabular}}
\end{table}

Notably:
(i) The fold-to-fold standard deviation \emph{tightens} on v3 for the two best-performing models (RF: 1.82 $\to$ 1.47; XGBoost: 1.61 $\to$ 1.26), exactly the qualitative effect one would expect if the cleaning protocol is removing label-noise outliers.
(ii) The RF vs.\ Chemprop $p$-value moves from $0.77$ on v2 to $0.16$ on v3---closer to but still well above the $\alpha = 0.05$ threshold. The shift is driven by a single Chemprop fold (fold~4, RMSE = 39.21~nm vs.\ the other four folds clustered at 30.5--33.9), which appears to be a borderline early-stopping outlier; the Chemprop std on v3 (3.27) is comparable to its v2 std (3.18), so the underlying variance is not protocol-driven.
(iii) The solvent-encoding ablation (the Effect of Solvent Information section of the main text, reported on v2) is expected to be preserved under v3 because the 1.81\% record loss affects both the +Solvent and Solute-only arms in the same way. A back-of-the-envelope spot check confirms this: holding the v2 Solute-only baselines fixed and pairing them with the v3 +Solvent numbers yields percent-RMSE reductions of 7.7\% (RF), 1.9\% (Chemprop), and 7.2\% (BiGRU); the Chemprop number is depressed by the same fold-4 outlier above, and we expect a full v3 retraining of both arms to recover the $\sim$6\% Chemprop reduction reported in Table~3.
(iv) Cross-dataset results (the Cross-Dataset Benchmarking section of the main text, Deep4Chem and Jung 2024) are unaffected by the primary-dataset protocol change because they train and evaluate on the external datasets' own train/val/test splits.

The full per-fold values for v3 are reported in Tables~\ref{tab:perfold_rf}--\ref{tab:perfold_chemberta}.


\section{Chromophore-Solvent Leakage Analysis}
\label{si:chromophore_leakage}

We audited the v3 5-fold stratified cross-validation splits used in main-text Table~2 for the chromophore-overlap concern raised by Greenman et al.~\cite{greenman2022multifidelity}: because the Joung+Beard dataset contains many more records (18{,}415) than unique solutes (7{,}846), the same chromophore can appear in both the training and test slices of a single fold under different solvents, giving the model an easier prediction target than would a scaffold-based split. For each fold and each model we partitioned the test records into two subsets---``chromophore seen in training'' (the same canonical solute SMILES appears in the train slice under any solvent) vs.\ ``chromophore novel'' (the solute is absent from train)---and recomputed RMSE/MAE/$R^2$ on each subset using the saved per-fold predictions. The fingerprint-based models (RF, XGBoost) were evaluated against the train-only slice they actually fit on; the deep learning models (BiGRU, Chemprop, ChemBERTa) were evaluated against the train+val slice they have access to (DL models use val for early stopping).

The overall leakage rate is high: 66.8\% of test records under the train-only partition and 69.3\% under the train+val partition; equivalently, 54.6\% of the unique solutes in any given test fold also appear in the corresponding train fold (Table~\ref{tab:si_leakage_perfold}).

\begin{table}[htbp]
\centering
\caption{Per-fold chromophore leakage statistics on the v3 5-fold stratified splits. ``Leaked vs.\ train'' counts test records whose canonical solute appears in the 64\% train slice (used by RF, XGBoost); ``leaked vs.\ train+val'' counts test records whose solute appears in the 80\% train+val slice (relevant for DL models with early stopping on val).}
\label{tab:si_leakage_perfold}
\small
\begin{tabular}{ccccc}
\toprule
\textbf{Fold} & \textbf{Test records} & \textbf{Unique test solutes} & \textbf{Leaked vs.\ train} & \textbf{Leaked vs.\ train+val} \\
\midrule
0 & 3683 & 2609 & 2440 (66.2\%) & 2534 (68.8\%) \\
1 & 3683 & 2564 & 2458 (66.7\%) & 2554 (69.3\%) \\
2 & 3683 & 2567 & 2469 (67.0\%) & 2546 (69.1\%) \\
3 & 3683 & 2547 & 2490 (67.6\%) & 2575 (69.9\%) \\
4 & 3683 & 2539 & 2451 (66.5\%) & 2545 (69.1\%) \\
\midrule
\textbf{Mean} & 3683 & 2565 & \textbf{2462 (66.8\%)} & \textbf{2551 (69.3\%)} \\
\bottomrule
\end{tabular}
\end{table}

Table~\ref{tab:si_leakage_per_model} reports the decomposition of each model's test-fold RMSE into chromophore-seen and chromophore-novel subsets. The chromophore-seen subset gives RMSEs in the 16--45~nm range, while the novel-solute subset gives RMSEs in the 50--70~nm range; the $\Delta$RMSE (novel minus seen) is $+25$ to $+36$~nm across all five model families.

\begin{table}[htbp]
\centering
\caption{Per-model RMSE on the chromophore-seen vs.\ chromophore-novel test subsets (mean $\pm$ std across five v3 folds). The full-test RMSE column reproduces the column from main-text Table~2 for reference. $\Delta$RMSE is the gap between novel and seen subsets.}
\label{tab:si_leakage_per_model}
\small
\resizebox{\textwidth}{!}{%
\begin{tabular}{lccccc}
\toprule
\textbf{Model} & \textbf{Leak rate} & \textbf{RMSE (seen)} & \textbf{RMSE (novel)} & $\boldsymbol{\Delta}$\textbf{RMSE} & \textbf{Full-test RMSE} \\
\midrule
RF (tuned) & 66.8\% & 16.36 $\pm$ 0.88 & 49.53 $\pm$ 2.90 & $+33.16$ & 31.50 $\pm$ 1.47 \\
XGBoost (MSE) & 66.8\% & 21.33 $\pm$ 0.29 & 50.52 $\pm$ 2.79 & $+29.19$ & 33.92 $\pm$ 1.26 \\
Chemprop (D-MPNN) & 69.3\% & 19.65 $\pm$ 4.10 & 51.90 $\pm$ 4.73 & $+32.26$ & 33.15 $\pm$ 3.27 \\
BiGRU (tuned) & 69.3\% & 20.96 $\pm$ 1.45 & 57.16 $\pm$ 3.03 & $+36.21$ & 36.20 $\pm$ 1.20 \\
ChemBERTa & 69.3\% & 45.13 $\pm$ 4.56 & 70.17 $\pm$ 3.68 & $+25.04$ & 54.09 $\pm$ 3.57 \\
\bottomrule
\end{tabular}}
\end{table}

Three observations:
(i) \textbf{All five model families benefit comparably from chromophore overlap.} $\Delta$RMSE spans a narrow 25--36~nm band across architectures spanning four families (tree ensemble, GNN, RNN, Transformer). No model is differentially advantaged: in particular, RF and Chemprop have $\Delta$RMSE of $+33.16$ and $+32.26$~nm respectively. The RF vs.\ Chemprop paired $t$-test on the chromophore-novel subset alone (n~=~5 folds) gives $t = -2.51$, $p = 0.066$---still failing to reject the null at $\alpha = 0.05$---so the RF$\approx$Chemprop statistical tie reported in main-text Table~2 ($p = 0.16$ on the full test set) is preserved when the comparison is restricted to chromophores genuinely absent from the training data.
(ii) \textbf{The novel-solute RMSEs roughly align with the wetlab OOD numbers.} The wetlab validation set (the Experimental Validation section of the main text, 16 commercial UV-absorbing compounds) gives RF MAE = 38.5~nm, Chemprop MAE = 26.2~nm, ChemBERTa MAE = 26.3~nm; the chromophore-novel subsets here give correspondingly elevated MAEs of comparable order of magnitude (full per-model MAE values in \texttt{results/r1\_9\_leakage\_per\_model.csv}). The pattern is consistent: in-distribution interpolation is substantially easier than true novel-chromophore prediction for all architectures, with the relative ranking of architectures depending on whether the test distribution is close to the training distribution or genuinely out of it.
(iii) \textbf{The magnitude of the effect is intrinsic to the dataset shape.} With $\sim$2.35 records per unique solute on average, any random or solvent-stratified split that does not group by solute will produce high chromophore-level overlap by construction. A scaffold-based split (our wetlab evaluation) is the methodological complement that controls for this. We retain the standard stratified-CV protocol in Table~2 because it is the convention used by Joung et al., Beard et al., Greenman et al., and Jung et al.\ on these same data, and matching the convention is necessary for apples-to-apples comparison with their published baselines; but readers should interpret those numbers as upper bounds on novel-chromophore generalisation.

The full analysis is reproducible from the saved v3 per-fold predictions via \texttt{python3 analyze\_chromophore\_leakage.py}, which writes \texttt{results/r1\_9\_leakage\_\{stats.json,per\_model.csv,per\_fold.csv,summary.md\}}.


\section{BiGRU Character-Type Saliency}
\label{sec:chartype_saliency}

Figure~\ref{fig:bigru_chartype} shows the mean gradient saliency aggregated by SMILES character type across all wetlab validation molecules.
Bond symbols (\texttt{=}, \texttt{\#}) receive the highest mean saliency, reflecting the critical role of double and triple bonds in establishing conjugated $\pi$-systems that determine UV absorption.
Heteroatoms (N, O, S) rank second, consistent with their role as electron donors and acceptors in chromophores, while aromatic atom characters and aliphatic carbons receive lower saliency, and brackets/syntax characters contribute minimally.

\begin{figure}[htbp]
 \centering
 \includegraphics[width=0.65\textwidth]{results/bigru_saliency_chartype.png}
 \caption{Mean BiGRU gradient saliency by SMILES character type, averaged over all wetlab validation molecules. Bond symbols (\texttt{=}, \texttt{\#}) and heteroatoms receive the highest saliency, consistent with the role of conjugated $\pi$-systems and electron-donating/withdrawing groups in determining UV absorption.}
 \label{fig:bigru_chartype}
\end{figure}


\section{Detailed Model Architectures}
\label{si:architectures}

This section collects the mathematical definitions and component-level explanations of the five model families benchmarked in the main text. The content was moved from the main-text Methods section to the Supporting Information on the grounds that the underlying architectures are documented in detail by their original publications (cited per subsection) and that the core findings of the manuscript---the locality argument, the RF$\approx$Chemprop tie, and the complementary OOD strengths of deep learning---do not depend on the equation-level details below.

\subsection*{Random Forest (RF)}

Random Forest~\cite{breiman2001random} is an ensemble method that constructs $B$ decision trees, each trained on a bootstrap sample of the data, and averages their predictions to reduce variance. For a regression task, the ensemble prediction is:
\begin{equation}
 \hat{y}(\vect{x}) = \frac{1}{B} \sum_{b=1}^{B} T_b(\vect{x})
 \label{eq:si_rf_predict}
\end{equation}
where $T_b$ denotes the $b$-th decision tree. Each tree partitions the feature space by recursively selecting the feature $j$ and threshold $s$ that minimise the mean squared error (MSE) within each resulting partition:
\begin{equation}
 (j^*, s^*) = \arg\min_{j, s} \left[ \sum_{x_i \in R_1(j,s)} (y_i - \bar{y}_{R_1})^2 + \sum_{x_i \in R_2(j,s)} (y_i - \bar{y}_{R_2})^2 \right]
 \label{eq:si_rf_split}
\end{equation}
where $R_1$ and $R_2$ are the two regions defined by the split and $\bar{y}_{R_k}$ is the mean target in region $R_k$. At each split, only a random subset of features (out of $d$ total) is considered, decorrelating the trees and improving generalisation.

The input representation for RF in this work is the Morgan circular fingerprint~\cite{rogers2010extended}, which encodes local chemical environments by enumerating all atom-centred circular substructures up to a specified radius $r$. Each substructure is hashed to a fixed-length binary vector of $n$ bits:
\begin{equation}
 \vect{x} = \text{MorganFP}(\text{mol},\, r=2,\, n=2048) \in \{0, 1\}^{2048}
 \label{eq:si_morgan_fp}
\end{equation}

\subsection*{XGBoost}

XGBoost~\cite{chen2016xgboost} constructs an additive ensemble of decision trees trained sequentially, where each new tree $f_t$ fits the residual errors of the current ensemble. The prediction after $T$ boosting rounds is:
\begin{equation}
 \hat{y}(\vect{x}) = \sum_{t=1}^{T} \eta \cdot f_t(\vect{x})
 \label{eq:si_xgb_predict}
\end{equation}
where $\eta$ is the learning rate (shrinkage parameter) that controls the contribution of each tree. Each tree $f_t$ is trained to minimise a regularised objective combining the loss on residuals and a penalty on tree complexity:
\begin{equation}
 \mathcal{L}^{(t)} = \sum_{i=1}^{n} l\!\left(y_i,\; \hat{y}_i^{(t-1)} + f_t(\vect{x}_i)\right) + \Omega(f_t)
 \label{eq:si_xgb_loss}
\end{equation}
where $\Omega(f_t) = \gamma \cdot |\text{leaves}| + \tfrac{1}{2}\lambda \sum_j w_j^2$ penalises the number of leaves and the magnitude of leaf weights $w_j$.

\subsection*{Recurrent Neural Networks and BiGRU}

Recurrent neural networks (RNNs)~\cite{elman1990finding} process sequential input by maintaining a hidden state $\vect{h}_t$ that is updated at each time step $t$:
\begin{equation}
 \vect{h}_t = \sigma(\matr{W}_h \vect{h}_{t-1} + \matr{W}_x \vect{x}_t + \vect{b})
 \label{eq:si_rnn}
\end{equation}
where $\vect{x}_t$ is the input at step $t$, $\matr{W}_h$ and $\matr{W}_x$ are weight matrices, and $\sigma$ is a nonlinear activation. Standard RNNs suffer from vanishing gradients over long sequences, motivating gated architectures.

The Gated Recurrent Unit (GRU)~\cite{cho2014learning} addresses this through two gating mechanisms, an update gate $\vect{z}_t$ and a reset gate $\vect{r}_t$, that control information flow:
\begin{align}
 \vect{z}_t &= \sigma(\matr{W}_z \vect{x}_t + \matr{U}_z \vect{h}_{t-1} + \vect{b}_z) \label{eq:si_gru_update} \\
 \vect{r}_t &= \sigma(\matr{W}_r \vect{x}_t + \matr{U}_r \vect{h}_{t-1} + \vect{b}_r) \label{eq:si_gru_reset} \\
 \tilde{\vect{h}}_t &= \tanh(\matr{W}_h \vect{x}_t + \matr{U}_h (\vect{r}_t \odot \vect{h}_{t-1}) + \vect{b}_h) \label{eq:si_gru_candidate} \\
 \vect{h}_t &= (1 - \vect{z}_t) \odot \vect{h}_{t-1} + \vect{z}_t \odot \tilde{\vect{h}}_t \label{eq:si_gru_output}
\end{align}
where $\sigma$ is the sigmoid function and $\odot$ denotes element-wise multiplication. The update gate $\vect{z}_t$ determines how much of the previous state to retain versus replace with the candidate $\tilde{\vect{h}}_t$, while the reset gate $\vect{r}_t$ controls how much past information is used to compute the candidate. This gating mechanism enables GRUs to capture long-range dependencies without the vanishing-gradient problem.

A Bidirectional GRU (BiGRU)~\cite{schuster1997bidirectional} processes the sequence in both directions, concatenating the forward and backward hidden states to capture context from both ends of the molecular string:
\begin{equation}
 \vect{h}_t^{\text{bi}} = [\overrightarrow{\vect{h}_t} \,\|\, \overleftarrow{\vect{h}_t}]
 \label{eq:si_bigru_concat}
\end{equation}

\subsection*{Graph Neural Network (D-MPNN)}

The Directed Message Passing Neural Network (D-MPNN)~\cite{yangAnalyzingLearnedMolecular2019} operates on the molecular graph by passing messages along directed bonds rather than atoms, avoiding redundant information flow to the source atom at each step. Given a molecular graph with atom features $\vect{x}_v$ and bond features $\vect{e}_{vw}$, directed bond hidden states $\vect{h}_{vw}^{(t)}$ are updated over $T$ steps:
\begin{equation}
 \vect{h}_{vw}^{(t)} = \tau\!\left(\vect{h}_{vw}^{(0)} + \matr{W}_t \sum_{u \in N(v) \setminus w} \vect{h}_{uv}^{(t-1)}\right)
 \label{eq:si_dmpnn}
\end{equation}
where $\tau$ is ReLU, $N(v) \setminus w$ denotes neighbours of $v$ excluding $w$, and $\vect{h}_{vw}^{(0)} = \tau(\matr{W}_0 \, [\vect{x}_v \,\|\, \vect{e}_{vw}])$ is initialised from concatenated atom and bond features. After $T$ message-passing steps, atom-level representations are obtained by summing incoming bond messages and passed through mean pooling to produce a fixed-length molecular vector.


\section{Chemprop Hyperparameter Probe}
\label{si:chemprop_hpo}

We ran an Optuna probe of Chemprop's architecture and regularisation on fold 0 of the v3 cleaned Joung+Beard dataset using a comparable 25-trial budget. The probe used the Tree-structured Parzen Estimator (TPE) sampler with $n_{\text{startup\_trials}} = 5$ and a MedianPruner ($n_{\text{warmup\_steps}} = 20$) to terminate clearly-underperforming runs early. The published default configuration was enqueued explicitly as the first trial so that every study contains an in-distribution baseline measurement.

\paragraph{Search space.}
The probe varied six hyperparameters with categorical distributions:
\begin{itemize}[leftmargin=2em,itemsep=2pt]
 \item \textbf{MPNN hidden dimension} $d_h \in \{200,\, 300,\, 400,\, 600,\, 900\}$ (default: 300)
 \item \textbf{Message-passing depth} $T \in \{3,\, 4,\, 5,\, 6\}$ (default: 3)
 \item \textbf{Dropout} $\in \{0.0,\, 0.1,\, 0.2,\, 0.3\}$ (applied to both BondMessagePassing and RegressionFFN; default: 0.0)
 \item \textbf{FFN depth} $\in \{1,\, 2,\, 3\}$ (default: 1)
 \item \textbf{FFN hidden dimension} $\in \{300,\, 500,\, 800\}$ (default: 300)
 \item \textbf{Batch size} $\in \{32,\, 50,\, 64,\, 128\}$ (default: 64)
\end{itemize}
The Noam learning-rate schedule (warmup $10^{-4} \to 10^{-3}$, exponential decay to $10^{-4}$), the maximum-epoch budget (100), the early-stopping patience (15 on val loss), and Chemprop's loss function were all held fixed at the values used for the default-config training, so the only variable across trials is the architecture/regularisation choice.

\paragraph{Best configuration.}
After 24 completed trials, the best fold-0 validation RMSE was achieved at trial 23 with
\[
 d_h = 900,\quad T = 4,\quad \text{dropout} = 0.1,\quad
 \text{FFN}\ 2\text{L} \times 800,\quad \text{batch} = 32.
\]
This configuration has 5.6M trainable parameters, about $9\times$ the 636K of the published default, and was re-trained on all 5 v3 folds with best-checkpoint restoration before computing test metrics; results are reported in main-text Table~2 as the ``Chemprop (tuned)'' row. Three additional independent runs of the same configuration during the probe gave fold-0 val RMSE values within $\pm 1$~nm of trial 23, indicating that the result is reproducible despite training stochasticity.

\paragraph{Per-fold comparison: tuned vs.\ default Chemprop on v3.}
The full 5-fold retrain with best-weight restoration (Table~\ref{tab:si_chemprop_tuned_perfold}) shows that the tuned configuration improves the cross-validated mean by $4.32$~nm and reduces fold-to-fold standard deviation from $3.27$ to $1.32$~nm. The improvement is driven primarily by fold 4, where the default-config Chemprop run produced an early-stopping outlier (RMSE $= 39.21$~nm) that the tuned configuration eliminates (RMSE $= 28.46$~nm, a $10.75$~nm drop).

\begin{table}[htbp]
\centering
\caption{Per-fold test-set RMSE (nm) for default vs.\ tuned Chemprop on the v3 cleaned Joung+Beard dataset, with best-checkpoint restoration on the validation set. Paired-$t$ and Wilcoxon signed-rank statistics computed across the five folds.}
\label{tab:si_chemprop_tuned_perfold}
\small
\begin{tabular}{cccc}
\toprule
\textbf{Fold} & \textbf{Default RMSE} & \textbf{Tuned RMSE} & $\boldsymbol{\Delta}$ \textbf{(tuned $-$ default)} \\
\midrule
0 & 30.62 & 30.70 & $+0.08$ \\
1 & 30.53 & 28.22 & $-2.31$ \\
2 & 31.49 & 26.92 & $-4.57$ \\
3 & 33.90 & 29.84 & $-4.06$ \\
4 & 39.21 & 28.46 & $-10.75$ \\
\midrule
\textbf{Mean $\pm$ Std} & \textbf{33.15 $\pm$ 3.27} & \textbf{28.83 $\pm$ 1.32} & \textbf{$-4.32$} \\
\bottomrule
\multicolumn{4}{l}{\footnotesize Paired $t$-test: $t = -2.40$, $p = 0.074$;\quad Wilcoxon signed-rank: $p = 0.125$.} \\
\end{tabular}
\end{table}

\paragraph{Per-fold comparison: tuned Chemprop vs.\ RF on v3.}
The headline comparison for the manuscript narrative is whether the tuned Chemprop dislodges RF as the leading model. The per-fold paired comparison gives a 2.67~nm mean advantage to tuned Chemprop, statistically significant under the parametric paired $t$-test ($t = -2.90$, $p = 0.044$) but \emph{not} significant under the conservative non-parametric Wilcoxon signed-rank test ($p = 0.125$, the minimum achievable $p$-value for $n = 5$ folds is $0.0625$; four of the five paired differences favour tuned Chemprop, the fifth favours RF by $0.37$~nm).

\begin{table}[htbp]
\centering
\caption{Per-fold test-set RMSE (nm) for tuned Chemprop vs.\ RF on the v3 cleaned Joung+Beard dataset.}
\label{tab:si_chemprop_tuned_vs_rf}
\small
\begin{tabular}{cccc}
\toprule
\textbf{Fold} & \textbf{RF RMSE} & \textbf{Tuned Chemprop RMSE} & $\boldsymbol{\Delta}$ \textbf{(tuned $-$ RF)} \\
\midrule
0 & 30.33 & 30.70 & $+0.37$ \\
1 & 30.46 & 28.22 & $-2.24$ \\
2 & 30.24 & 26.92 & $-3.32$ \\
3 & 32.67 & 29.84 & $-2.83$ \\
4 & 33.81 & 28.46 & $-5.35$ \\
\midrule
\textbf{Mean $\pm$ Std} & \textbf{31.50 $\pm$ 1.47} & \textbf{28.83 $\pm$ 1.32} & \textbf{$-2.67$} \\
\bottomrule
\multicolumn{4}{l}{\footnotesize Paired $t$-test: $t = -2.90$, $p = 0.044$;\quad Wilcoxon signed-rank: $p = 0.125$.} \\
\end{tabular}
\end{table}

\paragraph{Chromophore-novel subset.}
The same paired comparison restricted to test molecules whose canonical solute SMILES is absent from the corresponding train slice (the chromophore-novel subset defined in Section~\ref{si:chromophore_leakage}) gives tuned $= 46.54 \pm 2.49$~nm vs.\ default $= 51.90 \pm 4.23$~nm; mean $\Delta = -5.36$~nm, paired $t = -2.65$, $p = 0.057$ (Wilcoxon $p = 0.063$). The HPO benefit extends to genuinely novel chromophores, although the absolute novel-subset RMSEs remain $\sim$20~nm above the full-test numbers for both configurations.

\paragraph{Discussion.}
Three observations frame the manuscript treatment:
(i) The cross-validated RMSE improvement (4.32~nm) and the reduction of Chemprop's high default variance (3.27 $\to$ 1.32~nm) together indicate that Chemprop's published defaults are \emph{undertrained} for this dataset rather than near-optimal. The 5.6M-parameter tuned configuration corresponds to a Chemprop architecture that is 9$\times$ wider than the default; for the $\sim$12K training records per fold this larger model still benefits from training, contradicting the prior expectation that defaults would be production-validated across moderate-size molecular property datasets.
(ii) The tuned-vs-RF advantage is borderline under the conservative non-parametric test. We therefore retain the main-text headline language of a statistical tie between RF and Chemprop (as supported by both the v3 default comparison at $p = 0.16$ and the tuned-vs-RF Wilcoxon at $p = 0.125$) and report the tuned numbers as a supplementary row in Table~2 for transparency.
(iii) The default configuration is retained as the primary Table~2 comparison because it matches the evaluation convention used by published Chemprop benchmarks on these data~\cite{joung2021deep,beard2019comparative,greenman2022multifidelity,jung2024machine,chempropv2_2024}; matching the convention is necessary for apples-to-apples comparison with cited baselines, and the Chemprop v2 release paper explicitly recommends the defaults as production-ready out of the box for moderate-size molecular property tasks~\cite{chempropv2_2024}.

\paragraph{Reproducibility.}
The probe is reproducible from \texttt{tune\_chemprop.py} and the saved Optuna SQLite store \texttt{results/r1\_12\_chemprop\_optuna.db}; per-trial logs are in \texttt{logs/r1\_12\_chemprop\_trial\_*.log}. The 5-fold retrain with the best configuration is reproducible via \texttt{retrain\_chemprop\_tuned.py}, which calls \texttt{run\_chemprop.py} once per fold with the tuned hyperparameters injected through environment variables (\texttt{CHEMPROP\_D\_H}, \texttt{CHEMPROP\_DEPTH}, \texttt{CHEMPROP\_DROPOUT}, \texttt{CHEMPROP\_FFN\_N\_LAYERS}, \texttt{CHEMPROP\_FFN\_HIDDEN\_DIM}, \texttt{CHEMPROP\_BATCH\_SIZE}); per-fold predictions land in \texttt{results/chemprop\_tuned\_v3\_fold\{0..4\}\_*.npy} alongside the default-configuration files.


\section{Wetlab Chemical-Space Audit}
\label{si:wetlab_audit}

We computed two complementary chemical-space-overlap measures for every wetlab molecule against the full $7{,}846$-solute training set: (i)~nearest-neighbour Tanimoto similarity using Morgan fingerprints of radius~$2$ and $2048$ bits, and (ii)~exact-match check on the generic and specific Bemis--Murcko scaffolds. Per-molecule details are in Table~\ref{tab:si_wetlab_per_molecule}; the aggregate decomposition is in the four-bin Tanimoto histogram and the scaffold-match counts immediately below it.

\paragraph{Aggregate findings.}
At the SMILES level the wetlab set is novel: only Ferulic Acid is an exact duplicate of a training record (Tanimoto $= 1.000$), and the other 15 molecules are absent from the training data with a nearest-neighbour Tanimoto distribution of mean~$= 0.58$, median~$= 0.58$, range $[0.350, 0.676]$. The four-bin breakdown is:
\begin{itemize}[leftmargin=2em,itemsep=2pt]
 \item Near-duplicate (Tanimoto $> 0.85$): \textbf{1/16} (Ferulic Acid only).
 \item High-similar ($0.7 < $ Tanimoto $\leq 0.85$): \textbf{0/16}.
 \item Some-similar ($0.5 < $ Tanimoto $\leq 0.7$): \textbf{10/16}.
 \item Dissimilar (Tanimoto $\leq 0.5$): \textbf{5/16} --- Homosalate ($0.350$), Avobenzone ($0.379$), Amiloxate ($0.435$), Octocrylene ($0.458$), Octisalate ($0.491$).
\end{itemize}

\noindent
At the Bemis--Murcko scaffold level, however, the wetlab set is largely \emph{in-distribution}: \textbf{15/16} generic Murcko scaffolds and \textbf{15/16} specific Murcko scaffolds are present in the v3 training set. Only Avobenzone (a 1,3-diketone bridging two substituted phenyl rings) has a generic scaffold absent from training. The training set contains $3{,}013$ unique generic and $4{,}054$ unique specific Murcko scaffolds across its $7{,}846$ solutes, so the chromophore-class coverage is wide; the wetlab compounds therefore represent novel substitution patterns on chromophore scaffolds the models have seen at training time.

\paragraph{Strict-5 subset analysis (Tanimoto $\leq 0.5$).}
We additionally re-computed the per-model wetlab MAE on the five truly fingerprint-novel molecules to test whether the aggregate ranking is driven by the eleven scaffold-known compounds. Table~\ref{tab:si_wetlab_strict5} shows that the DL-vs-RF gap actually \emph{widens} on the strict-5 subset: the three DL models cluster at MAE $28$--$33$~nm while RF degrades to $46.1$~nm. This is consistent with the manuscript's complementary-strengths narrative: the DL architectures retain their out-of-distribution advantage over RF (and the ranking BiGRU $\sim$ Chemprop $<$ ChemBERTa $\ll$ RF is preserved) precisely on the molecules most chemically removed from the training data.

\begin{table}[htbp]
\centering
\caption{Per-model wetlab MAE (nm) on three nested subsets defined by Tanimoto similarity to the v3 training set: the full aggregate (16 molecules), the scaffold-known 11 (excluding Ferulic Acid and the five with Tanimoto $\leq 0.5$), and the strict-5 subset of truly fingerprint-novel molecules. Each subset row aggregates over the EtOH and MeOH measurements.}
\label{tab:si_wetlab_strict5}
\small
\resizebox{\textwidth}{!}{%
\begin{tabular}{lccccc}
\toprule
\textbf{Subset} & \textbf{BiGRU (def)} & \textbf{BiGRU (tuned)} & \textbf{RF} & \textbf{ChemBERTa} & \textbf{Chemprop} \\
\midrule
All 16 (aggregate) & 28.64 & 31.52 & 38.46 & 26.29 & 26.18 \\
Scaffold-known 11 (Tan $\in (0.5, 1)$) & 30.40 & 32.19 & 36.43 & 25.33 & 25.89 \\
Strict-5 (Tan $\leq 0.5$) & \textbf{28.35} & 35.66 & \textbf{46.09} & 33.00 & 28.46 \\
\bottomrule
\multicolumn{6}{l}{\footnotesize Strict-5 set: Amiloxate, Avobenzone, Homosalate, Octisalate, Octocrylene. Bold values denote} \\
\multicolumn{6}{l}{\footnotesize the strict-5 best (BiGRU default) and worst (RF). Numbers from \texttt{results/r1\_13\_subset\_mae.json}.} \\
\end{tabular}%
}
\end{table}

The single most striking strict-5 result is Octocrylene, where RF over-predicts by $87.7$~nm (and Chemprop by $81.3$~nm) while BiGRU's error is $42.5$~nm; the molecule's $\alpha$-cyano-cinnamate substitution and long branched ester chain shift its absorption blue of what fingerprint-based similarity to other cinnamates would predict. Conversely, on Amiloxate the D-MPNN's near-exact prediction (error $2.1$~nm) versus RF's $55.6$~nm over-prediction indicates that the multi-component graph encoder captures the long aliphatic ester linker as essentially absorption-inactive while RF's Morgan fingerprint treats the long-chain ester as a novel substructure.

\paragraph{Interpretation.}
Taken together, the audit places the wetlab benchmark in an honest middle ground rather than the genuine-OOD framing implied by the original ``15 entirely absent'' phrasing. The set is novel at the molecule level (substitution patterns, side chains, exact SMILES) but on chromophore scaffolds the models have seen in training. This is the operational deployment scenario most relevant to the photosafety- and UV-filter-screening applications the manuscript targets---chemists screening new ester/ether substitutions on known UV-absorber scaffolds---and is therefore a meaningful but specifically-scoped test. The closer approximation to scaffold-class extrapolation in the present manuscript is the Jung 2024 external benchmark ($\sim$50\% solute overlap with our primary training data, distinct experimental provenance, separate chromophore-class coverage); a dedicated scaffold-split evaluation across the primary v3 dataset is a priority for future work.

\paragraph{Reproducibility.}
The full audit is reproducible from the single script \texttt{analyze\_wetlab\_similarity.py} (Morgan-FP nearest-neighbour search + Murcko scaffold matching against the full v3 dedup CSV). The strict-5 subset metrics are reproducible from \texttt{results/wetlab\_predictions.csv} (per-molecule predictions for all five model families saved during the wetlab evaluation) and the helper output \texttt{results/r1\_13\_subset\_mae.json}; the per-molecule audit table and Tanimoto bin counts are written to \texttt{results/r1\_13\_wetlab\_per\_molecule.csv}, \texttt{results/r1\_13\_aggregate.json}, and \texttt{results/r1\_13\_summary.md}.

\begin{table}[htbp]
\centering
\caption{Per-molecule chemical-space audit of the 16 wetlab compounds against the v3 cleaned Joung+Beard training set ($18{,}415$ records, $7{,}846$ unique solutes). \textbf{Exact dup} = identical canonical SMILES present in train; \textbf{Gen-scaf} / \textbf{Spec-scaf} = generic / specific Bemis--Murcko scaffold present in train; \textbf{Nearest Tan} = Tanimoto similarity (Morgan radius~$2$, $2048$~bits) to nearest training solute; \textbf{Nearest train $\lambda_{\max}$} = mean experimental $\lambda_{\max}$ (nm) of that nearest training molecule. Full table with nearest-train SMILES and per-solvent experimental values in \texttt{results/r1\_13\_wetlab\_per\_molecule.csv}.}
\label{tab:si_wetlab_per_molecule}
\small
\resizebox{\textwidth}{!}{%
\begin{tabular}{lccccc}
\toprule
\textbf{Molecule} & \textbf{Exact dup} & \textbf{Gen-scaf} & \textbf{Spec-scaf} & \textbf{Nearest Tan} & \textbf{Nearest train $\lambda_{\max}$ (nm)} \\
\midrule
Amiloxate & N & Y & Y & 0.435 & 362.5 \\
Avobenzone & N & N & N & 0.379 & 286.2 \\
Caffeic Acid & N & Y & Y & 0.636 & 287.0 \\
Cinnamic Acid & N & Y & Y & 0.571 & 355.6 \\
Coumaric Acid & N & Y & Y & 0.586 & 355.6 \\
Dioxybenzone & N & Y & Y & 0.618 & 295.0 \\
Ferulic Acid & Y & Y & Y & 1.000 & 287.0 \\
Homosalate & N & Y & Y & 0.350 & 295.0 \\
Octinoxate & N & Y & Y & 0.667 & 322.5 \\
Octisalate & N & Y & Y & 0.491 & 316.0 \\
Octocrylene & N & Y & Y & 0.458 & 316.0 \\
Oxybenzone & N & Y & Y & 0.600 & 295.0 \\
PABA & N & Y & Y & 0.591 & 329.4 \\
Padimate O & N & Y & Y & 0.575 & 303.7 \\
Sinapic Acid & N & Y & Y & 0.676 & 287.0 \\
Sulisobenzone & N & Y & Y & 0.564 & 295.0 \\
\bottomrule
\end{tabular}}
\end{table}


\section{Solvent-Encoding Gain by Chromophore Multi-Solvent Coverage}
\label{si:multisolvent_audit}

We re-computed the +Solvent vs.\ Solute-only paired ablation restricted to test records whose canonical solute SMILES appears in the corresponding (train+val) slice with at least two distinct canonical solvent SMILES. The same per-fold trained checkpoints used for the main-text Table~3 are used here, so the only change is the subsetting; no new training was required.

\paragraph{Subset definitions.} For each fold $f$:
\begin{itemize}[leftmargin=2em,itemsep=2pt]
 \item Let $C^{\text{multi}}_f$ be the set of canonical solute SMILES that appear in the fold-$f$ train+val slice with $\geq 2$ distinct canonical solvent SMILES. Across the five folds, $|C^{\text{multi}}_f|$ ranges from $2{,}187$ to $2{,}208$ unique chromophores (about $28$\% of the $\sim$$7{,}846$ unique training solutes).
 \item The \textbf{multi-solvent test subset} is the set of test records whose canonical solute belongs to $C^{\text{multi}}_f$. Across folds this is $59.7$--$61.2$\% of test records (mean $= 60.3$\%).
 \item The \textbf{single-or-novel test subset} is the complement: test records whose canonical solute either appears in train+val with only one solvent or does not appear in train+val at all. This is the regime where solvent encoding has no per-chromophore signal to learn from.
\end{itemize}

\paragraph{Aggregate result (5-fold mean).} The full per-model decomposition is in Table~\ref{tab:si_multisolvent}. The headline observation is that the solvent-encoding contribution is much larger on the multi-solvent subset than the aggregate Table~3 numbers suggest:
\begin{itemize}[leftmargin=2em,itemsep=2pt]
 \item RF (tuned): aggregate gain $= 2.79$~nm (8.2\%); multi-solvent subset gain $= \mathbf{7.22}$~nm ($\mathbf{31.8}$\%); single-or-novel subset gain $= 0.42$~nm ($0.9$\%).
 \item Chemprop: aggregate gain $= 2.09$~nm (6.2\%); multi-solvent subset gain $= \mathbf{5.91}$~nm ($\mathbf{25.2}$\%); single-or-novel subset gain $= -0.28$~nm ($-0.6$\%).
 \item BiGRU (+Solv default): aggregate gain $= 2.58$~nm (6.6\%); multi-solvent subset gain $= \mathbf{5.65}$~nm ($\mathbf{21.2}$\%); single-or-novel subset gain $= 0.77$~nm ($1.5$\%).
\end{itemize}

\begin{table}[htbp]
\centering
\caption{Per-model decomposition of the solvent-encoding ablation by chromophore multi-solvent coverage. Test records are partitioned by whether the canonical solute appears in the corresponding (train+val) slice with $\geq 2$ distinct solvents (multi-solvent subset, $\sim$60\% of test) or with $\leq 1$ solvent / absent entirely (single-or-novel subset, $\sim$40\%). All values are 5-fold means computed from the same per-fold checkpoints used to construct Table~3. $\Delta$ denotes (Solute-only RMSE) $-$ (+Solvent RMSE); positive $\Delta$ means solvent encoding helps.}
\label{tab:si_multisolvent}
\small
\setlength{\tabcolsep}{4pt}
\resizebox{\textwidth}{!}{%
\begin{tabular}{lcccccccc}
\toprule
& \multicolumn{2}{c}{\textbf{Full (3{,}750)}} & \multicolumn{3}{c}{\textbf{Multi-solv subset ($\sim$2{,}260)}} & \multicolumn{3}{c}{\textbf{Single/novel subset ($\sim$1{,}490)}} \\
\cmidrule(lr){2-3} \cmidrule(lr){4-6} \cmidrule(lr){7-9}
\textbf{Model} & \textbf{+Solv} & \textbf{Solute} & \textbf{+Solv} & \textbf{Solute} & $\boldsymbol{\Delta}$ (\%) & \textbf{+Solv} & \textbf{Solute} & $\boldsymbol{\Delta}$ (\%) \\
\midrule
RF (tuned) & 31.34 & 34.13 & \textbf{15.49} & \textbf{22.72} & \textbf{7.22 (31.8\%)} & 45.85 & 46.27 & 0.42 (0.9\%) \\
Chemprop & 31.69 & 33.78 & \textbf{17.55} & \textbf{23.46} & \textbf{5.91 (25.2\%)} & 45.28 & 45.00 & $-0.28$ ($-0.6$\%) \\
BiGRU & 36.45 & 39.03 & \textbf{20.97} & \textbf{26.62} & \textbf{5.65 (21.2\%)} & 51.70 & 52.47 & 0.77 (1.5\%) \\
\bottomrule
\multicolumn{9}{l}{\footnotesize All RMSEs in nm. Analysis performed on v2 first-occurrence-dedup splits because both ablation arms (+Solvent and} \\
\multicolumn{9}{l}{\footnotesize Solute-only) were trained on v2; the v3 Greenman--Song cleaning removes records uniformly across chromophore-coverage} \\
\multicolumn{9}{l}{\footnotesize strata (Section~\ref{si:v2_v3_comparison} confirms the 1.81\% loss is independent of multi-solvent membership), so the multi-solvent} \\
\multicolumn{9}{l}{\footnotesize vs.\ single-or-novel ratios are expected to be preserved on v3.} \\
\end{tabular}%
}
\end{table}

\paragraph{Discussion.}
Three observations frame the manuscript treatment:
(i) Solvent encoding can only contribute predictive value where solvent-specific signal exists. For a chromophore the model has seen in only one solvent during training, the model has no basis to attribute the prediction's variability to the solvent feature versus the solute feature; the solvent vector becomes either an unused dimension or, at worst, an additional input the model has to learn to ignore. The multi-solvent subset is exactly the regime where this signal exists and the model can therefore exploit it.
(ii) The cross-architecture consistency reported in the main-text aggregate (6--8\%) is preserved and amplified on the multi-solvent subset (21--32\% across all three families), so the qualitative finding that concatenation-based solvent encoding helps regardless of architecture is unchanged.
(iii) The headline 21--32\% multi-solvent improvement matches the order-of-magnitude solvatochromic shifts reported in the photochemistry literature~\cite{reichardt2011solvents}, suggesting that the architectures are recovering close to the maximum extractable solvent-specific signal on this subset. The dilution to 6--8\% in the aggregate is a property of the train/test split's chromophore coverage, not a property of the architectures or the solvent-encoding strategy.

A consequence relevant to deployment: practitioners screening compounds whose chromophore class is well-represented in their training data (with multiple solvent observations) should expect solvent-encoding gains of $\sim$20--30\%; those screening novel chromophore classes should not expect solvent encoding to contribute much beyond a generic regularisation effect. The aggregate 6--8\% number is the right summary only for evaluation populations that resemble the cross-validated test mix.

\paragraph{Reproducibility.} The audit is reproducible from the single script \texttt{analyze\_multisolvent\_subset.py}, which loads the per-fold +Solvent and Solute-only prediction arrays from \texttt{results/} (saved during the existing v2 ablation runs), reconstructs the multi-solvent / single-or-novel partition from the v2 dataset (\texttt{previous\_code/UV\_canonical\_full\_dataset.csv}), and writes the per-fold and aggregate metrics to \texttt{results/r1\_14\_multisolvent\_subset.json} and \texttt{results/r1\_14\_summary.md}.


\section{Mamede Photosafety Dataset Reproducibility}
\label{si:mamede_reproducibility}

This section documents the full pipeline from upstream source to the dataset used for the Table~6 classifier training.

\paragraph{Upstream source.}
The dataset originates with Mamede, de Almeida, et al., \emph{Sci.\ Rep.}\ \textbf{2021}, \textit{11}, 3622 (DOI \texttt{10.1038/s41598-021-83104-4})~\cite{mamede2021machine}, who curated a set of $74{,}784$ compounds from Reaxys with binary photosafety labels (POS = absorbance in $290$--$700$~nm at MEC $\geq 1000$~L\,mol$^{-1}$\,cm$^{-1}$; NEG otherwise). The compound identifiers and SMILES strings were distributed in the article's Supplementary Information, which is publicly accessible without an institutional Reaxys subscription via the Springer Nature article page.

\paragraph{Local pipeline.}
\begin{enumerate}[leftmargin=2em,itemsep=2pt]
 \item \textbf{Download.} The Mamede SI table was downloaded as a tab-separated file ($74{,}784$ rows) containing three columns: \texttt{REAXYS\_Registry\_Number}, \texttt{SMILES}, and a Reaxys URL pointing back to the entry. The local copy is stored at \texttt{data/raw/mamede\_supplementary.txt} in the public repository.
 \item \textbf{Parse + canonicalise.} The pipeline script \texttt{postprocess\_reaxys.py} parses the SI file, drops the URL column, and passes each SMILES string through RDKit's \texttt{Chem.MolFromSmiles}\, $\rightarrow$\, \texttt{Chem.MolToSmiles(canonical=True)} round-trip. Of the $74{,}784$ entries, $74{,}783$ canonicalise successfully and $1$ SMILES fails to parse (logged but discarded), giving an effective recovery rate of $\sim$100\%. The output is \texttt{data/mamede\_classification\_dataset.csv}.
 \item \textbf{Label join.} Mamede's POS/NEG labels are present in the same SI file and are joined to the canonical-SMILES records without further filtering.
 \item \textbf{Train/test split.} The orchestrator script \texttt{run\_mamede.py rf\_cls} performs a single-level stratified split via \texttt{sklearn.model\_selection.train\_test\_split} with \texttt{test\_size=1996}, \texttt{stratify=class\_bin}, and \texttt{random\_state=7}, yielding $72{,}787$ train records, $998$ \texttt{test\_I} records, and $998$ \texttt{test\_II} records (the inner $998$/$998$ split uses a second \texttt{train\_test\_split} on the held-out 1996, also with \texttt{random\_state=7}). The Table~6 ``Our RF'' row aggregates predictions across \texttt{test\_I $\cup$ test\_II} ($N_{\text{test}} = 1{,}996$) to match the publication-scale evaluation.
 \item \textbf{Model.} RandomForestClassifier with $n_{\text{estimators}} = 500$, Morgan fingerprints of radius $2$ and $1024$ bits, default scikit-learn hyperparameters otherwise; same RNG seed (7). Out-of-bag score is reported in the per-run metrics file.
\end{enumerate}

\paragraph{Software versions.}
The pipeline uses RDKit $\geq$ 2024.03 (the project's pinned minimum in \texttt{requirements.txt}), scikit-learn 1.1+, and pandas/NumPy at their respective project pins. The full environment specification is in \texttt{pyproject.toml} and \texttt{requirements.txt} in the public repository.

\paragraph{Output artefacts.}
The per-run results file \texttt{results/mamede\_rf\_cls\_metrics.json} contains the aggregate \texttt{classification} block (TP/FP/TN/FN, sensitivity, specificity, accuracy, $F_1$, MCC, AUC), separate \texttt{test\_I} and \texttt{test\_II} blocks, the OOB score, and the configuration snapshot ($n_{\text{total}} = 74{,}783$, $n_{\text{train}} = 72{,}787$, $n_{\text{test1}} = 998$, $n_{\text{test2}} = 998$, $n_{\text{invalid\_smiles}} = 0$, seed = 7). All reported Table~6 numbers (sensitivity $= 0.876$, specificity $= 0.882$, accuracy $= 0.879$, AUC $= 0.942$) are read directly from this JSON.

\paragraph{Reproducibility command.}
From the public repository root:
\begin{verbatim}
python3 postprocess_reaxys.py # parse + canonicalise SI -> CSV
python3 run_mamede.py --model rf_cls # train RF + write metrics JSON
\end{verbatim}
Both scripts are pinned at \texttt{git} HEAD on the public repository (URL in the Data and Code Availability section).


\section{Experimental $\lambda_{\max}$ Uncertainty Audit}
\label{si:experimental_uncertainty}

To contextualise the magnitude of the model errors reported in Table~2, we quantified two complementary estimates of the inherent measurement noise in the source data. Both audits operate on the v2 first-occurrence-dedup combination of Joung 2020 and Beard 2019 ($18{,}755$ records pre-cleaning) because the v3 Greenman--Song protocol already removes duplicates with spread $> 5$~nm; the v2 view therefore retains the full noise distribution that the cleaning rule was designed to address.

\paragraph{Within-source duplicate spread.}
For each (canonical solute SMILES, canonical solvent SMILES) key with multiplicity $> 1$, we computed the within-key spread $\max(\lambda_{\max}) - \min(\lambda_{\max})$. Table~\ref{tab:si_uncertainty_within} reports the resulting distribution by source.

\begin{table}[htbp]
\centering
\caption{Within-source duplicate spread distribution. A multi-record key is a (canon, solvent) pair with at least two reported $\lambda_{\max}$ values in the named source. Median spread of $0$~nm indicates that more than half of all multi-record keys contain perfectly-agreeing values (typically re-extractions of the same primary measurement across multiple secondary citations). The $> 50$~nm tail captures large disagreements likely attributable to lab-to-lab protocol differences or labelling errors. The Greenman--Song cleaning rule (adopted in this revision per ) discards groups whose spread exceeds 5~nm; the ``$\leq 5$~nm'' fraction therefore quantifies the proportion of multi-record keys retained after cleaning.}
\label{tab:si_uncertainty_within}
\small
\begin{tabular}{lcccccc}
\toprule
\textbf{Source} & \textbf{Multi-keys} & \textbf{Median} & \textbf{p95} & \textbf{$\leq 5$~nm} & \textbf{$\leq 20$~nm} & \textbf{$> 50$~nm} \\
\midrule
Joung 2020 & 0 & -- & -- & -- & -- & -- \\
Beard 2019 & 273 & 0~nm & 21~nm & $89.7$\% & $94.1$\% & $1.1$\% \\
v2 combined & 294 & 0~nm & 38~nm & $86.1$\% & $90.5$\% & $4.1$\% \\
\bottomrule
\multicolumn{7}{l}{\footnotesize Joung 2020 reports each compound at most once per solvent in its published table, so it} \\
\multicolumn{7}{l}{\footnotesize contributes zero within-source multi-records. The v2-combined view picks up additional} \\
\multicolumn{7}{l}{\footnotesize multi-records from the Joung--Beard cross-source overlap (next paragraph).} \\
\end{tabular}
\end{table}

\paragraph{Inter-source disagreement (Joung vs.\ Beard).}
Of the $1{,}479$ (canon, solvent) keys in Joung 2020 and $16{,}999$ keys in Beard 2019, only $21$ keys are present in both. The cross-source disagreement on these $21$ paired observations is substantial: median $|\Delta\lambda_{\max}| = 49$~nm; mean = $36$~nm; p95 = $76$~nm; max = $77$~nm. Inter-source RMSD = $48$~nm. Only $38$\% of the paired observations agree within $5$~nm. The small sample size ($n = 21$) is the principal limitation here, and the paired compounds are biased toward chromophores of multi-paper solvatochromic interest, which are exactly the molecules most sensitive to subtle solvent-purity, temperature, and concentration differences across laboratories; the inter-source RMSD should therefore be treated as an upper bound on typical inter-laboratory disagreement rather than as a representative population statistic.

\paragraph{Implication for model error.}
The reported Table~2 RMSEs range from $31$~nm (RF) to $54$~nm (ChemBERTa). Comparing these to the within-source noise distribution: the model RMSEs are larger than the median or mean inter-duplicate spread but comparable to the $95$th percentile of within-source spreads ($21$~nm in Beard, $38$~nm in v2-combined). On the chromophores where the underlying measurements perfectly agree (the median-$=0$ regime), the model error is largely model-attributable; on the chromophores where the source data disagrees by tens of nm (the tail), the model error includes a substantial data-noise component. The headline RF$\approx$Chemprop tie reported in the main text is therefore consistent with both architectures operating near the data noise floor on a substantial fraction of test molecules---the regime in which further architectural improvement yields diminishing returns and where genuine error reduction requires either improved measurement protocols at data-curation time or models that explicitly carry an uncertainty estimate alongside the point prediction.

\paragraph{Reproducibility.}
The audit is reproducible from the single script \texttt{analyze\_experimental\_uncertainty.py}, which loads the cleaned per-source CSVs (\texttt{previous\_code/UV\_canonical\_dataset1.csv} for Joung and \texttt{previous\_code/UV\_canonical\_dataset2.csv} for Beard) plus the combined v2 file, and writes \texttt{results/r2\_1\_uncertainty.json} and \texttt{results/r2\_1\_uncertainty\_summary.md}.


\section{Per-Solute Solvent-Induced Shifts}
\label{si:solvent_shifts}

To quantify the empirical magnitude of solvatochromism in the v3 cleaned dataset directly from the data (without invoking polarity scales; see Section~\ref{si:polarity_breakdown} for the polarity-grouped analysis), we computed for each chromophore appearing under $\geq 2$ canonical solvents the span $\max(\lambda_{\max}) - \min(\lambda_{\max})$ across its observed solvent set. After collapsing intra-solvent duplicates by averaging, $2{,}516$ of the $7{,}846$ unique chromophores ($32$\%) satisfy the multi-solvent criterion. Figure~\ref{fig:si_solvent_shift_hist} shows the per-chromophore span distribution.

\begin{figure}[htbp]
\centering
\includegraphics[width=0.85\textwidth]{results/r2_4_solvent_shifts.png}
\caption{Distribution of per-chromophore solvent-induced spans on the v3 cleaned Joung+Beard dataset. The histogram is heavily right-skewed: median $= 10$~nm, $95$th percentile $= 52$~nm. The distribution tail is truncated at $200$~nm for visualisation; five chromophores have spans exceeding $200$~nm (the largest is Reichardt's $E_T(30)$ dye at $494$~nm, evaluated across $334$ different solvents in the dataset).}
\label{fig:si_solvent_shift_hist}
\end{figure}

\paragraph{Summary statistics.}
Median per-chromophore span = $10$~nm; mean = $16.3$~nm; p75 = $20$~nm; p90 = $36$~nm; p95 = $52$~nm; p99 = $123$~nm; max = $494$~nm. By bins:
\begin{itemize}[leftmargin=2em,itemsep=2pt]
 \item Span $= 0$~nm (chromophore reported with identical $\lambda_{\max}$ under multiple solvents): $4.8$\%.
 \item Span $\leq 5$~nm (within instrument precision): $26.0$\%.
 \item Span $\leq 10$~nm (typical solvatochromism): $51.6$\%.
 \item Span $\leq 20$~nm: $75.0$\%.
 \item Span $> 50$~nm (large solvatochromic shifts): $5.6$\%.
 \item Span $> 100$~nm (extreme shifts; merocyanines, push--pull cyanines, $E_T(30)$-class dyes): $1.4$\%.
\end{itemize}

\paragraph{Top-15 chromophores by span.}
Table~\ref{tab:si_solvent_shift_top} reports the fifteen largest per-chromophore spans, together with structural class flags from a lightweight SMARTS classifier. The list is dominated by the chromophore classes that the solvatochromism literature predicts: the $E_T(30)$ reference dye~\cite{reichardt2011solvents} at the top, push--pull merocyanines (intramolecular charge transfer between donor and acceptor through a conjugated polyene bridge), and stilbene-pyridinium cyanines (charged acceptor + amine donor across a stilbene bridge). The empirical ranking confirms that the v3 dataset captures the full operational range of solvatochromic behaviour relevant for UV/Vis prediction model evaluation.

\begin{table}[htbp]
\centering
\caption{Top-15 chromophores by solvent-induced span on the v3 cleaned dataset. ``n solv'' = number of distinct canonical solvents the chromophore appears under in v3. ``Span'' = $\max(\lambda_{\max}) - \min(\lambda_{\max})$ across those solvents. ``Class'' = lightweight SMARTS structural flags (empty when no rule matches; many strong-solvatochromism chromophores are push--pull merocyanines that do not match a simple SMARTS).}
\label{tab:si_solvent_shift_top}
\small
\setlength{\tabcolsep}{3pt}
\begin{tabular}{rrlp{0.55\textwidth}}
\toprule
\textbf{Span (nm)} & \textbf{n solv} & \textbf{Class} & \textbf{Canonical SMILES} \\
\midrule
493.5 & 334 & -- & Reichardt's $E_T(30)$ betaine (zwitterionic phenoxide-pyridinium) \\
418.0 & 2 & -- & \texttt{CC(=O)[O-]} (acetate, edge-case data point) \\
235.0 & 2 & stilbene & polythiophene-bridged porphyrin (stilbene-like) \\
225.0 & 9 & -- & merocyanine (bis-benzimidazole + thiabarbituric acid) \\
210.0 & 10 & stilbene & stilbene-pyridinium + hydroxynaphthalene \\
163.3 & 9 & stilbene & stilbene-quinolinium + coumarin oxide \\
153.0 & 3 & -- & merocyanine (cyclobutenedione + dimethylanilino) \\
153.0 & 12 & stilbene & stilbene-pyridinium + hydroxyphenyl \\
152.0 & 3 & -- & merocyanine (cyclobutenedione + anilino) \\
151.0 & 3 & -- & merocyanine (cyclobutenedione + dimethylanilino) \\
151.0 & 3 & -- & merocyanine (cyclobutenedione + dimethylanilino) \\
150.0 & 3 & -- & merocyanine (cyclobutenedione + dimethylanilino) \\
148.0 & 3 & -- & push--pull coumarin + indolium + sulfonate \\
145.0 & 13 & -- & extended-aromatic acceptor + naphthyl-aryl wings \\
145.0 & 3 & -- & push--pull coumarin + indolium + sulfonate \\
\bottomrule
\multicolumn{4}{l}{\footnotesize Full SMILES strings for all $2{,}516$ multi-solvent chromophores in} \\
\multicolumn{4}{l}{\footnotesize \texttt{results/r2\_4\_solvent\_shifts.csv}.} \\
\end{tabular}
\end{table}

\paragraph{Reproducibility.}
Reproducible from the single script \texttt{analyze\_solvent\_shifts.py}. Outputs: \texttt{results/r2\_4\_solvent\_shifts.json}, \texttt{results/r2\_4\_solvent\_shifts.csv} (full per-chromophore table), \texttt{results/r2\_4\_solvent\_shifts.png} (the histogram in Figure~\ref{fig:si_solvent_shift_hist}), and \texttt{results/r2\_4\_summary.md}.


\section{Solvent Polarity Stratification}
\label{si:polarity_breakdown}

To test whether any solvent polarity class is disproportionately difficult for the four headline models, we partitioned the $18{,}415$ v3 records into four polarity tiers based on standard $E_T(30)$-style groupings of canonical solvent SMILES (Reichardt 2011~\cite{reichardt2011solvents}), plus an ``other'' bucket for solvents not represented in the standard tier definitions:
\begin{itemize}[leftmargin=2em,itemsep=2pt]
 \item \textbf{Nonpolar} ($n = 3{,}010$, $16.4$\%): n-hexane, n-heptane, n-octane, cyclohexane, benzene, toluene, xylene, CCl$_4$.
 \item \textbf{Moderate} ($n = 6{,}619$, $35.9$\%): DCM, CHCl$_3$, 1,2-DCE, THF, 2-Me-THF, EtOAc, MeOAc, 1,4-dioxane, acetone.
 \item \textbf{Polar aprotic} ($n = 3{,}625$, $19.7$\%): MeCN, DMF, DMSO, propionitrile, DMA, GBL, NMP.
 \item \textbf{Polar protic} ($n = 3{,}412$, $18.5$\%): water, MeOH, EtOH, $n$-PrOH, $i$-PrOH, $n$-BuOH, $n$-PenOH, ethylene glycol, glycerol.
 \item \textbf{Other} ($n = 1{,}749$, $9.5$\%): long-tail solvents not in the standard tier definitions (mixed solvents, less-common substituted aromatics, etc.).
\end{itemize}

\paragraph{Test-fold RMSE by tier.}
Table~\ref{tab:si_polarity_breakdown} reports pooled-5-fold test RMSE per (model, tier) cell. Three observations frame the headline:

\begin{table}[htbp]
\centering
\caption{Pooled 5-fold test-set RMSE (nm) per solvent-polarity tier. ``$n$'' = pooled test-set size for that tier (sum across the five folds). ``ALL'' row aggregates the full v3 test set and matches the headline Table~2 numbers exactly.}
\label{tab:si_polarity_breakdown}
\small
\begin{tabular}{lrcccc}
\toprule
\textbf{Tier} & $\boldsymbol{n}$ & \textbf{RF} & \textbf{Chemprop} & \textbf{BiGRU} & \textbf{ChemBERTa} \\
\midrule
Nonpolar & 3{,}010 & \textbf{27.7} & 28.4 & 30.5 & 50.5 \\
Moderate & 6{,}619 & 34.5 & 37.9 & 39.6 & 53.6 \\
Polar aprotic & 3{,}625 & 29.7 & 32.0 & 32.6 & 42.5 \\
Polar protic & 3{,}412 & 30.2 & 29.7 & 32.7 & 49.7 \\
Other & 1{,}749 & 32.4 & 31.8 & 44.7 & \textbf{84.4} \\
\midrule
\textbf{ALL} & 18{,}415 & \textbf{31.5} & 33.3 & 36.2 & 54.2 \\
\bottomrule
\end{tabular}
\end{table}

\begin{itemize}[leftmargin=2em,itemsep=2pt]
 \item \textbf{Cross-architecture ranking is preserved on every tier.} RF $\lesssim$ Chemprop $<$ BiGRU $\ll$ ChemBERTa holds in nonpolar, moderate, polar aprotic, polar protic, and ``other'' tiers. The statistical-tie claim and the complementary-strengths narrative are not artefacts of a specific solvent class.
 \item \textbf{Nonpolar is the easiest tier across all four models.} RF reaches $27.7$~nm RMSE on nonpolar (vs.\ $31.5$~nm on ALL), Chemprop reaches $28.4$~nm. This is consistent with two compounding effects: (a) nonpolar test records skew toward less complex chromophores due to solubility constraints (highly conjugated push--pull dyes are typically insoluble in hexane/cyclohexane), and (b) absence of strong specific solute--solvent interactions (hydrogen bonding, dipole-dipole) reduces measurement variance.
 \item \textbf{The ``other'' tier reveals ChemBERTa's failure mode most starkly.} On the long-tail solvents not represented in the standard tier definitions, ChemBERTa's RMSE jumps to $84.4$~nm (vs.\ $54.2$~nm aggregate), while RF and Chemprop stay near their aggregate values ($32.4$ and $31.8$~nm). This is consistent with the prediction-ceiling artefact diagnosed in Section~5.2 of the main text and analysed in more detail in SI Section~\ref{si:chemberta_ceiling}: ChemBERTa's degradation mode is concentrated on inputs that drift far from its pretraining distribution, of which the ``other''-tier solvents are a representative sample.
\end{itemize}

There is no evidence of a polarity tier where one architecture systematically wins or loses against its peers; the relative ranking is preserved with $\pm 5$~nm shifts at most. The polarity-grouped analysis therefore reinforces rather than complicates the manuscript's headline cross-architecture findings.

\paragraph{Reproducibility.}
Reproducible from the single script \texttt{analyze\_solvent\_polarity.py}, which loads the per-fold v3 predictions for all four headline models, joins them to the v3 dataset's canonical solvent column, applies the tier mapping, and writes per-(model, tier) RMSE and MAE to \texttt{results/r2\_5\_polarity\_breakdown.\{json,csv,png\}} and \texttt{results/r2\_5\_summary.md}.


\section{ChemBERTa $\sim$556~nm Prediction-Ceiling Diagnostic}
\label{si:chemberta_ceiling}

We tested three hypotheses for the prediction-ceiling artefact visible in Figure~5, where ChemBERTa's predictions for long-wavelength chromophores cluster near $556$~nm regardless of the experimental $\lambda_{\max}$ value:

\paragraph{(H1) Training-distribution imbalance.}
Fold-0 train ($n = 11{,}785$) contains $980$ records ($8.3$\%) with $\lambda_{\max} > 600$~nm, $272$ ($2.3$\%) with $> 700$~nm, and only $47$ ($0.4$\%) with $> 800$~nm. The train median is $396$~nm, p95 = $647$~nm, p99 = $755$~nm. The deep-coloured chromophores ($\lambda_{\max} > 600$~nm) are a long thin tail in the training distribution: present but underrepresented relative to the bulk.

\paragraph{(H2) ChemBERTa output range.}
Across the full pooled-5-fold test set ($18{,}415$ records), ChemBERTa's predictions span $[185, 556.4]$~nm. The \emph{maximum} predicted value is $556.4$~nm; $0$\% of predictions exceed $600$~nm. On the $1{,}555$ test records with experimental $\lambda_{\max} > 600$~nm, the median ChemBERTa prediction is $549$~nm, $37$\% of predictions fall in the $[550, 650]$~nm band, $98.5$\% are underpredicted by more than $50$~nm, and the median underprediction is $126$~nm. This is a hard ceiling on the model's predictive range, not a gradual underprediction.

\paragraph{(H3) SMILES-length truncation.}
Combined solute--solvent SMILES strings have median length $61$ chars, p95 $135$ chars, max $645$ chars. Only $0.23$\% of test records exceed the MAX\_LEN $= 256$ truncation cap in the ChemBERTa tokenizer. Mean SMILES length does increase with $\lambda_{\max}$ (60.8 chars at $\lambda_{\max} < 400$~nm $\to$ 101.6 chars at $\lambda_{\max} > 700$~nm), but the truncation rate per bin is uniformly low (Table~\ref{tab:si_chemberta_smileslen}). H3 is therefore ruled out as the primary cause of the $\sim$556~nm ceiling.

\begin{table}[htbp]
\centering
\caption{Combined solute--solvent SMILES length statistics by experimental $\lambda_{\max}$ bin (v3 pooled-5-fold test set). ``Frac trunc'' = fraction of records in that bin whose combined SMILES exceeds the MAX\_LEN $= 256$ token cap used by the ChemBERTa tokenizer.}
\label{tab:si_chemberta_smileslen}
\small
\resizebox{\textwidth}{!}{%
\begin{tabular}{lrrr}
\toprule
\textbf{$\lambda_{\max}$ bin (nm)} & $\boldsymbol{n}$ & \textbf{Mean SMILES len (chars)} & \textbf{Frac trunc at MAX\_LEN $= 256$} \\
\midrule
$150$--$400$ & 9{,}490 & 60.8 & $0.18$\% \\
$400$--$500$ & 5{,}037 & 70.4 & $0.20$\% \\
$500$--$600$ & 2{,}318 & 84.8 & $0.47$\% \\
$600$--$700$ & 1{,}139 & 90.1 & $0.26$\% \\
$700$--$1100$ & 430 & 101.6 & $0.23$\% \\
\bottomrule
\end{tabular}%
}
\end{table}

\paragraph{Diagnosis.}
The $\sim$556~nm ceiling arises from the conjunction of (H1) and (H2), with (H3) ruled out. The proximate cause is the regression head's output range. The deeper cause is the pretrained-representation mismatch: ChemBERTa was pretrained on ZINC (\texttt{seyonec/ChemBERTa-zinc-base-v1}), whose drug-like molecules are dominated by short-wavelength absorbers; the [CLS] embeddings ChemBERTa produces for long-wavelength chromophores (extended polyenes, push--pull cyanines, BODIPYs, porphyrins) are not well-separated in the pretrained representation from those of medium-wavelength absorbers, and the regression head fine-tuned on top of these collapsed embeddings cannot push its predictions beyond the $\sim$556~nm boundary even when shown long-wavelength training targets, because the input representations are not sufficiently distinct. The training-target imbalance (H1) compounds this: with $\sim$92\% of training mass below $600$~nm, the MAE-minimising regression head defaults to predictions near the training median for any input it has not seen many times.

This is consistent with the broader finding in the deep-learning literature that fine-tuning a pretrained transformer cannot recover representation distinctions that the pretraining task did not require, even if downstream training data is provided~\cite{ross2022molformer}; the right remedy is either to pretrain on a UV-Vis-rich corpus (which would require a different pretraining task than ZINC's drug-likeness) or to augment the v3 training data with long-wavelength chromophores (which the field is already actively working on; see Future Work).

\paragraph{Three-panel diagnostic figure.}
Figure~\ref{fig:si_chemberta_diagnostic} shows (a) the train vs.\ test $\lambda_{\max}$ histograms with the 600~nm line marked, (b) the ChemBERTa parity plot colour-coded by whether the input SMILES would have been truncated at MAX\_LEN $= 256$ (truncated records are a small minority and are scattered, not concentrated at the ceiling), and (c) mean combined-SMILES length by $\lambda_{\max}$ bin with the MAX\_LEN cap marked.

\begin{figure}[htbp]
\centering
\includegraphics[width=\textwidth]{results/r1_16_chemberta_diagnostic.png}
\caption{Three-panel ChemBERTa prediction-ceiling diagnostic. (a) v3 fold-0 train + pooled-test $\lambda_{\max}$ distribution; only $8.3$\% of train records exceed 600~nm. (b) ChemBERTa parity plot, colour-coded by whether the combined solute--solvent SMILES exceeds MAX\_LEN $= 256$; the horizontal cluster at $\sim$556~nm is unrelated to truncation status (most truncated records are scattered, not concentrated at the ceiling). (c) Mean combined-SMILES length by $\lambda_{\max}$ bin; long-wavelength chromophores have longer SMILES but the MAX\_LEN cap is rarely hit (median length stays $\leq 100$ chars across all bins, with $\leq 0.5$\% truncation in any bin).}
\label{fig:si_chemberta_diagnostic}
\end{figure}

\paragraph{Reproducibility.}
Reproducible from the single script \texttt{analyze\_chemberta\_ceiling.py}, which loads the per-fold ChemBERTa predictions from \texttt{results/chemberta\_v3\_fold\{0..4\}\_*.npy} and the v3 dataset CSV, computes the three-hypothesis diagnostics, and writes \texttt{results/r1\_16\_chemberta\_diagnostic.\{json,png\}}, \texttt{results/r1\_16\_chemberta\_diagnostic\_top\_y.csv} (per-record diagnostic for the 200 test records with highest experimental $\lambda_{\max}$), and \texttt{results/r1\_16\_summary.md}.

\bibliography{Proposal}

\end{document}
